%% LyX 2.5.1 created this file.  For more info, see https://www.lyx.org/.
%% Do not edit unless you really know what you are doing.
\documentclass[brazil]{article}
\usepackage[T1]{fontenc}
\usepackage[utf8]{luainputenc}
\usepackage{babel}
\begin{document}

\section{Marxismo}

\subsection{Marxismos}

Como a base da minha análise da sociedade é o marxismo, precisamos começar por aqui. Começamos então com a questão mais óbvia possível: o que é o marxismo? A verdade é que não há um consenso sobre o que é marxismo. Há diversas vertentes que podem ser mencionadas, e cada vertente propõe uma própria definição sobre o que é marxismo e a partir desta definição, ela interpreta que a sua própria vertente é o ``verdadeiro'' marxismo.

Ingo Elbe\cite{elbe2021entre} vai discutir a existência de três marxismos que divergem por uma questão filosófica: Marxismo Tradicional, Marxismo Ocidental e Nova Leitura de Marx. Dentro da vertente tradicional ainda podemos encontrar uma divisão entre ``stalinistas'' e ``trotskistas'', que tem como um dos principais pontos de divergência a reivindicação das experiências socialistas reais e as implicações que decorrem disto. Se quisermos prosseguir com mais subdivisões, dentro do ``stalinismo'' (marxismo-leninismo) podemos encontrar mais um leque de diferentes vertentes como maoismo, hoxhaismo, etc.

Mesmo quando duas vertentes concordam em algum nível (exemplo: o que define marxismo é o método) elas podem discordar em outro nível (exemplo: qual exatamente é o método). Eu não defendo que exista nenhum ``verdadeiro'' marxismo no sentido mais tradicional. O que me interessa é identificar qual o melhor procedimento para interpretar e modificar a realidade.
\begin{quote}
Em matéria de marxismo, a ortodoxia se refere antes e exclusivamente a um método\cite{lukács1969historia}.
\end{quote}
Se entendermos marxismo como uma simples continuidade do trabalho de Marx, todas as vertentes que partem de Marx, em maior ou menor grau, se prestam a esse papel e tornam-se igualmente verdadeiras. 
\begin{quote}
A tradição de todas as gerações mortas oprime como um pesadelo o cérebro dos vivos. E justamente quando parecem empenhados em revolucionar-se a si e às coisas, em criar algo que jamais existiu, precisamente nesses períodos de crise revolucionária, os homens conjuram ansiosamente em seu auxilio os espíritos do passado, tomando-lhes emprestado os nomes, os gritos de guerra e as roupagens, a fim de apresentar e nessa linguagem emprestada\cite{marx1852brumario}.
\end{quote}
Agora se adicionarmos a exigência de que seja coerente com os escritos originais, precisamos refletir sobre mais uma coisa: Sequer todo o escrito de Marx é coerente entre si? E se não, a vertente deve ser coerente com que parte do trabalho de Marx? A pergunta é mais difícil do que parece. Mas eu tendo a pensar que precisamos reconhecer a existência de incoerências entre diferentes obras do Marx. Marx morreu com mais de 60 anos e por toda sua vida adulta estudou e produziu trabalhos intelectuais. É de se esperar que ao longo do tempo algumas ideias tenham sofrido alteração.

Ao longo da obra de Marx, existem ao menos seis exposições da forma-valor, sendo todas diferentes entre si\cite{barreira2022nova}. Ainda há uma mudança no foco de estudo de Marx que se reflete no foco das suas obras publicadas, partindo de obras mais puramente filosóficas como “A Sagrada Família” e se movendo na direção de obras que passam a ter a economia política como tema central, como o próprio “O Capital”. O quanto de filosofia persiste em suas obras posteriores é assunto de polêmicas. Esta discussão atinge tal nível que alguns autores distinguem em Marx a existência de um “Marx esotérico” e um “Marx exotérico”\cite{elbe2021entre}, outros discutem um corte epistemológico entre o “jovem Marx” (fortemente influenciado por Hegel e Feuerbach;) e o “velho Marx” (quando há uma ruptura de Marx com a influência desses dois filósofos)\cite{lira2021althusser}. E até mesmo o próprio Marx chega a ser acusado de simplificar suas formulações teóricas\cite{barreira2022nova}.

Um relato da trajetória de estudos do Marx é dado pelo próprio no prefácio de Para a Crítica da Economia Política.
\begin{quote}
O meu estudo universitário foi o da jurisprudência, o qual no entanto só prossegui como disciplina subordinada a par de filosofia e história. No ano de 1842-43, como redactor da Rheinische Zeitung, vi-me pela primeira vez, perplexo, perante a dificuldade de ter também de dizer alguma coisa sobre o que se designa por interesses materiais. Os debates do Landtag Renano sobre roubo de lenha e parcelamento da propriedade fundiária, a polémica oficial que Herr von Schaper, então Oberprásident da província renana, abriu com a Rheinische Zeitung sobre a situação dos camponeses do Mosela, por fim as discussões sobre livre-cambismo e tarifas alfandegárias proteccionistas deram-me os primeiros motivos para que me ocupasse com questões económicas\cite{marx2008critica}.
\end{quote}
Entendo que Marx possui um conjunto de obras que permite que você escolha diferentes recortes e desenvolva uma dada teoria que seja consistente com tal recorte escolhido. Este é um dos motivos por que não considero uma boa meta a busca tradicional pelo ``verdadeiro'' marxismo. De fato, não considero sequer existir um verdadeiro marxismo. Diversas vertentes buscam desenvolver diferentes momentos, ênfases e obras de Marx. Mas nenhuma é o próprio Marx continuando sua obra, logo. Defendo assim, que outro critério deve ser adotado.

Defendo a ideia de abandonar o personalismo e a própria figura do Marx para visar a formulação do marxismo que melhor descreva a realidade. Partindo de uma perspectiva científica, adoto a postura de que o marxismo deve ter mais compromisso com a adequada descrição do que pode ser observado, do que respeito com as ideias de qualquer indivíduo do passado. Evidentemente a correta descrição vai precisar de pressupostos corretos e é aqui que reside a herança de Marx. Entendo então, que por marxista, quero dizer que defendo que a correta descrição da realidade é obtida quando adotamos a concepção de mundo defendida por Marx e Engels, popularmente conhecida como materialismo dialético.

\subsection{Concepção de Mundo}

Por materialismo dialético eu quero dizer, em poucas palavras, que esta é a concepção de mundo que encontramos exposta na obra “Anti-Dühring”.
\begin{quote}
\ldots{} o ‘sistema’ do Senhor Dühring, de que este livro é uma crítica, estende-se a domínios teóricos muito vastos: tive de segui-lo por toda Parte e opor às suas concepções as minhas. Assim, a crítica negativa resultou positiva; a polêmica transformou-se em exposição mais ou menos coerente do método dialético e da ideologia comunista defendida por Marx e por mim, numa série de domínios bastante vastos \ldots{}

\ldots{} Uma observação de passagem: tendo sido criada por Marx, e em menor escala por mim, a concepção exposta neste livro, não conviria que eu a publicasse à revelia do meu amigo. Li-lhe o manuscrito inteiro antes da impressão; e o décimo capítulo da Parte segunda, consagrada à economia política (Sobre a história crítica) foi escrito por Marx. Infelizmente, eu o tive de resumir por motivos extrínsecos. Era, aliás, hábito nosso ajudarmo-nos mutuamente na especialização de cada um\cite{engels1877antiduhring}.
\end{quote}
Mais do que um conjunto de leis rígidas ou um método, materialismo dialético é para mim uma concepção de mundo pautada em alguns pressupostos gerais sobre como é este mundo.
\begin{quote}
Uma concepção de mundo abrange a totalidade das opiniões e concepções de um indivíduo ou sociedade sobre o mundo, sobre nós mesmos como seres humanos e sobre a vida e a posição dos seres humanos no mundo\cite{worldview2023}.
\end{quote}
É a partir desta obra que gosto de introduzir noções de materialismo dialético e dialética materialista.
\begin{quote}
O\textbf{ Materialismo Dialético} é uma compreensão científica da matéria, da consciência e da relação entre ambas. Ele é usado para entender o mundo através do estudo dessas relações.

A \textbf{Dialética Materialista} é uma ciência que estuda as leis gerais do movimento, da mudança e do desenvolvimento da natureza, da sociedade e do pensamento humano\cite{worldview2023}.
\end{quote}
Sob este viés, por materialismo dialético temos então uma uma concepção materialista do mundo que em sua totalidade, também incorpora uma relação dialética entre os entes que constitui esse mundo. Esta concepção de mundo, por ser materialista, implica ideia da existência de um mundo que existe objetivamente independente de qualquer consciência, portanto, anterior à mente, e consequentemente regido por leis naturais e impessoais.

Ao chamar a dialética materialista por ciência, não queremos que a palavra “ciência” aqui seja entendida no sentido moderno, como uma disciplina da universidade. Mas sim como um conjunto de conhecimentos sobre um dado assunto e uma área de interesse de investigação. Assim tenho por dialética materialista o conjunto de pressupostos (e a investigação) destes pressupostos que nos orientaram sobre como as entidades que compõem o mundo se relacionam e se transformam. Ela diz respeito sobre como o mundo evolui. A dialética materialista vai exigir uma concepção de mundo onde ele esteja em constante transformação e totalmente conectado, em outras palavras, não há ente no mundo nem isolado e nem que permaneça inalterado.

Materialismo dialético e dialética materialista formam então uma só teoria que deve também englobar tanto como concebemos o mundo tanto quanto sobre o nosso lugar nele. Não devem vê-los como duas coisas separadas e unidas mecanicamente: materialismo e dialética. Temos uma concepção de mundo materialista que engloba também uma perspectiva dialética de como este mundo se transforma.

Estas definições são bastante genéricas e flexíveis. Não são nada mais do que alguns pressupostos básicos sobre como a constituição e desenvolvimento da natureza. Por marxismo, ainda faço uma adição de um posicionamento político bem definido em defesa das massas. Isto é algo que não surge necessariamente apenas de uma correta compreensão do mundo, porém é uma parte intrínseca do movimento político que recebeu este título. Entendo então o marxismo tanto quanto um movimento científico em busca da adequada descrição da realidade pautado em pressupostos materialistas e dialéticos, quanto um movimento político de transformação dessa realidade em prol das massas, necessariamente.

Esta forma particular de conceber o marxismo não é livre de polêmicas. Há quase tantas concepções de dialética quanto há marxistas no mundo Como será explorado nas próximas sessões, estou consciente de que minha interpretação rompe com a tradição filosófica tradicional e hegeliana. De fato, eu particularmente tenho preferência por chamar o conjunto de ideias que defendo por ``teoria marxista'' do que ``filosofia marxista'', seguindo a observação feita por Nildo Viana:
\begin{quote}
A tarefa da filosofia (que depois será também alvo de crítica de Marx e será substituída por teoria, entre outros termos ocasionais, como “socialismo científico”, “ciência” e, após ele, será chamado marxismo ou teoria)...\cite{marx2020introducao}
\end{quote}
Nesta forma de conceber a teoria marxista, Hegel passa a ser visto como uma influência filosófica em Marx, mas não necessariamente um componente crucial do marxismo. O marxismo assim exposto, torna-se independente de Hegel e sua dialética.
\begin{quote}
Gramsci, nos Cadernos do Cárcere, nos diz que é necessário entender que a filosofia da práxis nasceu de “aforismos critérios práticos por um mero acaso” e nos explica que Marx dedicou sua força intelectual de forma sistemática a outros problemas, principalmente de natureza econômica. Sua concepção de mundo aparece então apenas de forma implícita nestes aforismos. 

Uma das causas do erro pelo qual se vai em busca de uma filosofia geral que esteja na base da filosofia da práxis e se nega implicitamente a esta uma originalidade de conteúdo e de método parece que reside no seguinte: que se faz confusão entre a cultura filosófica pessoal do fundador da filosofia da práxis, ou seja, entre as correntes filosóficas e os grandes filósofos pelos quais ele se interessou fortemente em sua juventude e cuja linguagem reproduz frequentemente\ldots e as origens ou as partes constitutivas da filosofia da práxis\ldots{} os elementos de spinozismo, de feuerbachismo, de hegelismo, de materialismo francês etc. não são de modo algum partes essenciais da filosofia da práxis\cite{gramsci2024cadernos}.
\end{quote}
Não tenho como negar que há uma influência hegeliana em Marx durante toda a sua vida, e que esta influência vai deixar marcas. Mas argumento que conforme à teoria marxista se desenvolve, já em Marx e Engels, essa influência vai sendo engolida por práticas que são mais científicas que filosóficas. Vejo ecos desta interpretação principalmente na obra de Engels.
\begin{quote}
Logo que descobrirmos -- e afinal de contas ninguém mais do que Hegel nos ajudou a descobri-lo -- que, assim colocada, a tarefa da filosofia se reduz a pretender que um filósofo isolado realize aquilo que somente a humanidade em seu conjunto poderá realizar, em seu desenvolvimento progressivo -- assim que descobrirmos isso a filosofia, no sentido tradicional da palavra, chega a seu fim. Já não interessa a “verdade absoluta”, inatingível por este caminho e inacessível ao único indivíduo, e o que se procura são as verdades relativas, adquiridas através das ciências positivas e da generalização de seus resultados por meio do pensamento dialético. A filosofia, em seu conjunto, termina com Hegel; por um lado, porque em seu sistema se resume, da maneira mais grandiosa, todo o desenvolvimento filosófico; por outro lado, porque este filósofo nos indica, ainda que inconscientemente, a saída desse labirinto dos sistemas para o conhecimento positivo e real do mundo.

\ldots{}

Esta interpretação põe fim à filosofia no campo da história, exatamente da mesma forma que a concepção dialética da natureza torna a filosofia da natureza tão desnecessária quanto impossível. Agora, já não se trata de tirar do cérebro as conexões entre as coisas, mas de descobri-las nos próprios fatos. Expulsa da natureza e da história, só resta à filosofia um único refúgio: o reino do pensamento puro, no que dele ainda está de pé: a doutrina das leis do próprio processo do pensamento, a lógica e a dialética\cite{engels1886feuerbach}.
\end{quote}
Mas também na obra de Marx.
\begin{quote}
É preciso “deixar a filosofia de lado”, é preciso desembarcar dela e dedicar-se como um homem comum ao estudo da realidade, tarefa para a qual existe uma gigantesca quantidade de material literário \ldots{} A relação entre filosofia e estudo do mundo real corresponde à relação entre onanismo e amor sexual\cite{marx2007ideologia}.
\end{quote}
É preciso comentar que apesar de ser uma posição não-tradicional, há outros marxistas que adotam uma posição próxima a minha. Paul Cockshott, em seu blog pessoal não apenas questiona o fato de uma fração dos marxistas verem como vital o estudo de Hegel para o marxismo, mas não terem a mesma opinião acerca de economistas como David Ricardo e Adam Smith (que também foram fortes influências em Marx), argumenta que vê no estudo de Hegel um potencial maior de confundir do que de esclarecer.
\begin{quote}
Taimur, você certamente é um palestrante muito bom e claro, mas não acho que a palestra tenha me tornado mais simpático à ideia de estudantes contemporâneos estudarem Hegel. Longe de ser uma ajuda, creio que seria um obstáculo ao seu desenvolvimento intelectual.

\ldots{}

Temos tantas ferramentas para analisar o mundo, desenvolvidas nos últimos 200 anos, que retroceder à década de 1820 seria um terrível passo retrógrado\cite{cockshott2020hegel}.
\end{quote}
A ideia central da crítica é relativamente simples: Hegel tem o mérito de ser um pensador avançado para o seu tempo e uma influência positiva em Marx dado o contexto no qual estava inserido. Contribuindo de forma inegável para o desenvolvimento do que hoje temos por marxismo. Mas seus métodos não constituem mais necessariamente a forma mais eficiente de entender a realidade. Pautando o objetivo do marxismo como sendo a melhor investigação e compreensão da realidade, resta ao marxista buscar as melhores ferramentas disponíveis para entender o mundo, e esta ferramenta não é, necessariamente, a dialética hegeliana.

\subsection{Um não-Hegeliano}

Devido a posição central que a dialética ocupa na polêmica que envolve Marx, Engels e Hegel, eu acredito ser necessario esclarecer que meu afastamento de Hegel se dá de forma consciente. É comum a ideia de que há um método em Marx no sentido hegeliano. A partir disso, algumas pessoas então vão se dedicar a procurar a dialética hegeliana oculta em o “O Capital”. Eu reconheço que em “O Capital” há uma tentativa por Marx de usar a dialética mais no sentido concebido por Hegel do que Engels, especialmente no primeiro capítulo, conforme confessado pelo próprio Marx. O que coloco em questão e tanto a profundidade quanto a relevância deste esforço. 
\begin{quote}
Confessei-me, portanto, abertamente discípulo daquele grande pensador e coqueteei mesmo aqui e ali no capítulo sobre a teoria do valor com o modo de expressão que lhe é peculiar\cite{marx2013capital1}.
\end{quote}
Os manuscritos do Grundrisse são possívelmente o mais perto que chegamos de um trabalho de Marx com este tom mais profundamente filosófico e hegeliano. Porém este mesmo projeto foi abandonado por Marx e o resultado disso é que torna-se necessário procurar o método dialético considerando que ele existe mas está oculto em suas obras publicadas, além de manuscritos abandonados. Já Engels, diferente de Marx, escreveu de forma mais propositiva sobre filosofia. Porém Engels se afasta mais de Hegel que Marx se afastou. É mais fácil encontrar elementos da dialética de Hegel no “O Capital” do que no “Anti-Dühring”. 

Porém é interessante recordarmos que esta obra foi revisada por Marx antes da publicação após ser decidido em comum acordo entre ambos sobre a necessidade dela e de que incumbiria a Engels escrevê-la, conforme relatado por José Paulo Netto em um texto introdutório para a edição da Boitempo do Anti-Dühring\cite{engels1877antiduhring}. Além disso Marx contribuiu diretamente para a obra escrevendo um dos capítulos a recomendou em cartas. Evidentemente isto não implica que ambos compartilhem da mesma percepção de dialética, mas me parece razoável assumir que a concepção engeliana não era considerada um absurdo para Marx.
\begin{quote}
\ldots{} e se o Sr. Most falhou em notar que há muito a ser aprendido com as exposições positivas de Engels {[}no “Anti-Dühring”{]}, não apenas por trabalhadores comuns e até mesmo ex-trabalhadores como ele, que se supõem capazes de conhecer tudo e se pronunciar sobre tudo no menor tempo possível, mas até mesmo por pessoas cientificamente educadas, então eu só posso ter pena dele por sua falta de julgamento\cite{marx1877bracke}.
\end{quote}
Não entendo que a concepção de dialética hegeliana que por ventura Marx possuísse seja sua principal contribuição, e consequentemente não é isto que vai definir o marxismo. No próprio ``O Capital'', entendo que temos mais uma organização da exposição das idéias seguindo um modelo típico dialético, do que descobertas que foram feitas devido a aplicação do método de investigação dialético conforme o a concepção hegeliana. O citado primeiro capítulo no qual Marx quis render homenagem ao Hegel é reescrito já na segunda edição pelo próprio Marx. Essa posição não é algo inteiramente nova, Adorno discutindo uma possível interpretação equivocada da dialética escreve:
\begin{quote}
se denomina dialética apenas um método determinado de pensamento, uma maneira determinada de se apresentar {[}darstellen{]} a coisa {[}Sache{]} - tal como Marx designou, numa ocasião talvez não exatamente feliz\cite{adorno2021introducao}
\end{quote}
(Adorno em Introdução à Dialética)

Esta vem acompanhada com a seguinte nota de rodapé:
\begin{quote}
Adorno se refere aqui manifestamente a uma passagem no posfácio de Karl Marx à segunda edição de O capital, a qual conduziu a uma torrente de interpretações conflitantes com respeito ao significado e alcance da dialética em Marx... Marx defende aqui seu " método dialético"{} contra seus críticos, diferenciando o " modo de apresentação"{} do " modo de investigação" . Isso conduziu à questão de se, segundo a compreensão marxiana, dialética seria apenas uma forma de apresentação científica de um assunto (Sache), ou se (também) seria a lei histórica e, correlativamente, genética da própria coisa apresentada\cite{adorno2021introducao}.
\end{quote}
Dessa forma, me parece que a dialética hegeliana no “O Capital”, acaba ganhando uma sobrevida por uma questão estética de como os resultados foram apresentados, não sendo igualmente claro que tenha se provado ser o método último de investigação da realidade. É popular encontrarmos nesta discussão a distinção entre método de exposição e método de investigação. Assim, o método de exposição apresentado no “O Capital” é de relativamente fácil observação pois é diretamente acessível pela própria exposição das ideias no livro, já quanto ao método de investigação os estudiosos sofrem para tentar “reconstruí-lo” investigando “método oculto” no livro.

Marx comenta a existência de um “método oculto” em “Para a Crítica da Economia Política”\cite{barreira2022nova}. Por isso eu entendo a existência desta discussão acerca da centralidade de Hegel em Marx. Minha crítica está na avaliação de que esta reconstrução seja produtiva. Prosseguindo a discussão, precisamos então discutir “qual” é a dialética que eu considero importante. Eu defendo que seja o conjunto de pressupostos gerais como apresentado por Engels, ou seja, princípios de como o mundo se transforma e evolui. Na minha visão, esta é uma formulação mais interessante da dialética. Dessa forma, ela funciona mais como um guia para a investigação, como um conjunto de princípios norteadores, do que um método em si. Além de ser também menos nebuloso e mais flexível.

Ao meu lado, tenho todos grandes movimentos políticos derivados da obra marxista que obteve êxito em algum grau, todas as revoluções socialistas que se inspiraram nas obras de Marx. Todos estes movimentos se se basearam em uma teoria no qual não teve especial carinho pela dialética hegeliana. O que acredito que levanta um debate saudável sobre quais são os aspectos da teoria marxista que através da prática se provaram válidos e fundamentais. 

Reforço então que tenho por marxista, uma questão de princípios gerais que moldam nossa concepção do mundo e sobre como nos posicionamos neste mundo. Isto é, defender os pressupostos materialistas de como o mundo é (matéria antes da mente, etc) em conjunto com os pressupostos dialéticos de como esse mundo material evolui (tudo está em constante transformação, etc), associado com uma questão política sobre como nos posicionamos e atuamos politicamente no mundo (em defesa da classe dos trabalhadores e pela superação do capitalismo).

\subsection{O que resta da filosofia}

Para fechar essa discussão, gostaria de dispensar a leitura mais tradicional do papel da filosofia. Praticamente desde que me interessei por marxismo eu tenho lido sobre a filosofia. Isso decorreu do fato de que eu sempre fui crítico e cético quanto à filosofia. Porém eu gosto de entender o que quero criticar, logo, senti a necessidade de entender melhor a filosofia

Uma das primeiras coisas que tive a oportunidade de aprender foi de que não há um consenso sobre o que é filosofia. Cada pensador que ganhou o direito de ser reconhecido como filósofo definiu a filosofia de uma forma diferente. E isso de certa forma fornece uma certa flexibilidade para que a filosofia seja defendida de forma abstrata. Em um momento filosofia é ``pensar de forma crítica'' (e ninguém vai se opor à esta tarefa), mas em outro momento não existe filosofia chinesa pois a filosofia é definida como uma tradição que nasceu na Grécia antiga. 

As duas alegações nunca podem serem defendidas em conjunto, pois implicaria em dizer que os chineses não eram capazes de pensar de forma crítica. Me parece que em grande parte a definição mais regular de filosofia na academia é aquela que é dada pela disciplina de história da filosofia colocando sua origem necessariamente no mundo grego antigo. É isto que vamos encontrar na coleção de “História Da Filosofia” de Giovanni Reale\cite{reale1990historia}.

Este projeto de história da filosofia que trata ela como um ``produto exclusivo do gênio grego'', me parece um projeto que nasceu intrinsecamente como um projeto eurocêntrico, desde Hegel. E normalmente, quando eu teço críticas a filosofia, minha crítica se refere principalmente a esta tradição com origem em platão. Outra descoberta que tive foi de que também não há também consenso sobre o que é a filosofia marxista. Porém, para que o materialismo dialético na formulação marxista-leninista recebesse o título de filosofia, foi necessário adotar uma definição de filosofia que fosse adequada. Pode ser lido em “Dialectical Materialism” do Spirkin\cite{spirkin1983dialectical}, por exemplo, que filosofia é a unidade de uma concepção de mundo e um método de investigação, parafraseando o próprio Engels quando descreve a teoria que desenvolveu junto de Marx como uma concepção de mundo comunista e método de investigação dialético\cite{engels1877antiduhring}. Esta é uma posição similar ao que podemos podemos encontrar em Gramsci:
\begin{quote}
É preciso,portanto, demonstrar preliminarmente que todos os homens são “filósofos”, definindo os limites e as características desta “filosofia espontânea”, própria de “todo o mundo”, isto é, da filosofia que está contida: 1) na própria linguagem, que é um conjunto de noções e de conceitos determinados e não, apenas, de palavras gramaticalmente vazias de conteúdo; 2) no senso comum e no bom senso; 3) na religião popular e, portanto, também em todo o sistema de crenças, superstições, opiniões, modos de ver e de agir que se descortinam naquilo que geralmente se chama de “folclore”\cite{gramsci2024cadernos}.
\end{quote}
O fato é que o materialismo dialético na concepção marxista leninista, é significativamente diferente da filosofia como é tradicionalmente concebida. Algo que o próprio Spirkin reconhece e o leva a classificar o marxismo como uma filosofia científica, e mais precisamente, como a primeira filosofia verdadeiramente científica. Um título que talvez não faça muito sentido da forma que a filosofia é tradicionalmente concebida. 

Essa ruptura com a filosofia clássica fica evidente por exemplo em Engels quando escreve o “Ludwig Feuerbach e o fim da filosofia clássica alemã” e o “Anti-Dühring”. Não por coincidência, meus textos preferidos sobre o assunto. Mas como o próprio Engels coloca, isso não implica que necessariamente tudo da filosofia clássica tenha se tornado obsoleto e que a própria tarefa de filosofar não possua mais nenhuma utilidade. Eu defendo apenas uma redução de escopo e ambição. É nesse sentido que aprecio o “Investigações Filosóficas” do Wittgenstein\cite{wittgenstein1999investigacoes}. 

Entendo hoje, que a filosofia, em sua formulação mais tradicional e abrangente, isto é, a filosofia que não é marxista, mas também não se restringe à tradição iniciada em Platão, mas sim uma reflexão especulativa geral sobre tudo, não deve ser uma teoria que visa explicar o mundo. Mas sim que a filosofia deve ser encarada como uma atividade de esclarecimento sobre o uso da linguagem: ela dissolve pseudoproblemas ao revelar quando estamos usando a linguagem além dos limites do sentido.

Me parece que Wittgenstein, de certa forma, assim como Engels, não elimina a filosofia por completo, mas altera a ideia sobre qual é o papel da filosofia. Ao invés de vermos na filosofia o papel de questionar como o mundo é, transferimos esta responsabilidade para as ciências e permitimos que a filosofia tenha por atividade questionar como usamos a linguagem. A filosofia se torna então uma prática de esclarecimento da linguagem. Eu tendo a pensar que a filosofia na prática, que eu não iria discordar da utilidade, é exatamente isso. Quando se pergunta " Deus existe?"{} A gente acaba tendo que discutir o que queremos dizer com " Deus" . E talvez até mesmo o que quer dizer com " existir" . Se tivermos tudo esclarecido (e nem tudo pode ser esclarecido), podemos então buscar a comprovação ou não de sua existência através da ciência. Esta é a discussão que entendo ser por natureza filosófica. 

Dessa cabe a ciência produzir conhecimento, e não a filosofia. A filosofia, pela forma que opera, pelos limites da lógica, não é capaz de produzir conhecimento, pois ela não se baseia na observação e extração de 'verdade' nenhuma do mundo real. Acredito que minha posição é coerente com as ideias expostas por Engels.
\begin{quote}
Por toda a parte, não se trata mais de congeminar conexões na cabeça, mas de as descobrir nos factos. Para a filosofia desalojada da Natureza e da história, fica ainda então apenas o reino do pensamento puro, na medida em que ainda resta: a doutrina das leis do próprio processo do pensar, a lógica e dialéctica\cite{engels1886feuerbach}.
\end{quote}
Agora que já esclarecemos como a filosofia encontrou seu auge e fim com Hegel e que já reduzi a filosofia clássica à função de resolver os problemas de linguagem, eu queria discutir a relação entre ciência e filosofia,e de certa forma, contestar a ideia de que a filosofia seja a mãe da ciência. De forma simples, meu argumento é de que a ciência atual e a filosofia atual ambas se originam em uma mesma tradição mais antiga, que pode ser chamada tanto como “filosofia universal” quanto como “ciência universal”. Isto é, não é nem a filosofia atual e nem a ciência atual.
\begin{quote}
Assim que compreendermos --- e, em definitivo, ninguém nos ajudou mais a essa compreensão do que o próprio Hegel --- que a tarefa da filosofia, colocada dessa maneira, não significa senão a tarefa de que um filósofo singular deve realizar aquilo que só a humanidade inteira no seu desenvolvimento progressivo pode realizar --- assim que compreendermos isto, estará também no fim toda a filosofia no sentido da palavra até aqui. Abandona-se a «verdade absoluta», inalcançável por esta via e por cada um individualmente, e, em troca, perseguimos as verdades relativas alcançáveis pela via das ciências positivas e do compêndio dos seus resultados por intermédio do pensar dialéctico. Com Hegel, rematà-se, em geral, a filosofia; por um lado, porque ele reuniu todo o desenvolvimento dela no seu sistema, da maneira mais grandiosa; por outro lado, porque, se bem que inconscientemente, ele nos mostra o caminho {[}que nos leva{]} deste labirinto dos sistemas ao conhecimento positivo real do mundo\cite{engels1886feuerbach}. 
\end{quote}
Além disso, a ciência moderna deve mais ao conhecimento prático e empírico acumulado por séculos pelas classes baixas, do que às abstrações frequentemente especulativas e erradas dos filósofos que pertenciam à classe dominante. Esse apagamento histórico de um lado e exaltação além da realidade do outro é, ao meu ver, fruto direto da luta de classes.
\begin{quote}
Basta dizer que, antes de 1600, antes do " nascimento da ciência e da filosofia modernas" , a colheita de frutos da árvore da filosofia era bastante escassa.

\ldots{}

Por outro lado, os antigos gregos e romanos e os antigos egípcios demonstraram uma arte e uma habilidade maravilhosas na construção e mesmo na engenharia mecânica. As suas abóbadas e arcos, que suportavam cargas tão pesadas, foram construídos com tal precisão e força que indicam um conhecimento altamente sofisticado da mecânica prática. Seja como for, o conhecimento destes antigos construtores e engenheiros não era científico, mas tecnológico. Era puramente empírico, não se baseava em princípios filosóficos\cite{frank1952origin}.
\end{quote}
Começamos então com estes dados históricos. Platão, que está na raiz da filosofia ocidental moderna, teria nascido uns 4 séculos antes de cristo, o que nos dá uma janela de tempo de 2 mil anos de existência da filosofia sem que desse muitos frutos. Se formos voltar a Tales de Mileto, precisamos adicionar mais 2 séculos de improdutividade da filosofia. Só isto já deveria ser o suficiente para questionar. E este período de escassa produção corresponde exatamente ao período anterior ao nascimento da ciência moderna. Isto é, a filosofia só foi capaz de produzir frutos quando a ciência moderna surgiu. Esta declaração revela uma posição semelhante à de Einstein:
\begin{quote}
Admiramos a Grécia antiga porque fez nascer a ciência ocidental\ldots{} Mas para atingir uma ciência que descreva a realidade, ainda faltava uma segunda base fundamental que, até Kepler e Galileu, foi ignorada por todos os filósofos. Porque o pensamento lógico, por si mesmo, não pode oferecer nenhum conhecimento tirado do mundo da experiência. Ora, todo o conhecimento da realidade vem da experiência e a ela se refere. Por este fato, conhecimentos, deduzidos por via puramente lógica, seriam diante da realidade estritamente vazios. Desse modo Galileu, graças ao conhecimento empírico, e sobretudo por ter se bat\cite{einstein1981comovejo}ido violentamente para impô-lo, tornou-se o pai da física moderna e provavelmente de todas as ciências da natureza em geral. Se, portanto, a experiência inaugura, descreve e propõe uma síntese da realidade...
\end{quote}
Acredito que esta analogia sobre\cite{spinoza1663prefacio} a colheita da árvore da filosofia venha da afirmação de Descartes de que a ‘utilidade’ da filosofia está nos seus frutos. Evidentemente isso não significa que as civilizações do passado não detinham conhecimento, pelo contrário. Porém, se pegarmos o conhecimento do mundo antigo, que era um conhecimento menos geral e mais técnico voltado ao desenvolvimento de tecnologias, isto existiu de forma independente da filosofia. Frequentemente se fala isso sobre outras nações como Índia e China, mas muitos tentam colocar civilizações ocidentais como ``superiores'' pois supostamente não tinham ``apenas tecnologia, tinham filosofia''. Porém, como podemos ver desde o surgimento da filosofia, até o nascimento da ciência, ela não produziu grandes frutos, nem mesmo no mundo ocidental.
\begin{quote}
O grande fosso que existiu durante a Antiguidade e a Idade Média entre " princípios"{} e " saber-fazer" , entre " ciência"{} e " tecnologia" , estava intimamente ligado à estrutura social. O " saber-fazer"{} do artesão e do engenheiro era adquirido através do trabalho físico e não do puro esforço intelectual. Era o conhecimento empírico do homem de baixo estatuto social, do artesão e até do escravo. O conhecimento filosófico era cultivado e promovido pelos homens de estatuto social elevado, nomeadamente os padres e os estadistas. Os estratos " inferiores"{} coleccionavam " factos" ; os superiores avançavam princípios. O contacto entre os dois tipos de conhecimento era desencorajado pelos costumes sociais. Se um homem de estatuto social elevado tentasse aplicar a sua " filosofia"{} ou " ciência"{} ao trabalho do artesão, era condenado ao ostracismo\cite{frank1952origin}.
\end{quote}
Temos então uma separação social, uma questão de classes na valorização do conhecimento produzido pelos trabalhadores de classes baixas e pelos intelectuais da classe alta. Ou seja, não é apenas que o conhecimento empírico não se baseia na filosofia, mas o conhecimento empírico estava fatalmente separado do conhecimento filosófico. Não apenas eram áreas distintas, como as pessoas que praticavam eram pessoas distintas. 

E ainda mais, eram pessoas que não tinham contato entre si pois pertenciam a diferentes classes. Se alguém de classe alta fosse ``pego'' trabalhando com um conhecimento empírico, seria punido socialmente. Então aqui a gente tem que a evolução destas diferentes áreas do conhecimento ocorreu de forma independente, por pessoas diferentes, de classes sociais diferentes. 

Eu não estou negando que em torno do século XV, quando surgiu o que hoje chamamos de ciência, este processo esteve relacionado com o que hoje chamamos de filosofia. Devido a mudanças sociais na época, membros da classe alta com mais tempo livre para abstrair questões práticas, treinados na filosofia, tiveram contato e se interessaram então por este conhecimento técnico, criando as condições necessárias para o surgimento de em uma ciência mais geral e abstrata. 

Se quisermos, podemos localizar aqui a contribuição da filosofia. Porém precebo que neste processo de transformação do conhecimento técnico em ciência, a separação que existia entre filosofia e conhecimento técnico, se transforma na separação entre filosofia e ciência. E é preciso notar que descrevendo a história da ciência desta forma, temos um conhecimento empírico que se transforma em ciência, beneficiada do contato com a filosofia, mas está mesmo permanece tanto separada do conhecimento técnico no passado, quanto da ciência no presente\footnote{Estou ciente que dizer que permanece separado é uma afirmação polêmica. Mas leia de forma bastante aberta tendo por mente que os cientistas modernos tradicionalmente não estudam filosofia. Isto é, produzem ciência sem ter formação e frequemtemente nem interesse em filosofia.}.

Ainda sim, a história do desenvolvimento da ciência moderna está atrelada também às necessidades do modo de produção moderno. Se reconhece que o conhecimento técnico nasceu em resposta às necessidades materiais, por exemplo, a matemática como necessidade de calcular áreas e controlar o imposto, mas a ciência moderna também responde a condições materiais específicas. O desenvolvimento da termodinâmica não é apenas causa da revolução industrial, mas também consequência.

Assimo desenvolvimento da ciência também responde a uma questão social. Porém, da mesma forma que a astronomia surge por aqueles que se interessavam em astrologia, ou a química por quem se interessava por alquimia, a ciência surge por aqueles que se interessavam por filosofia. Dessa filosofia pré-1600, considerada então uma espécie de ciência/filosofia universal, surge então a “filosofia própria” e a “ciência própria”. Poderíamos então dizer que a ciência é a mãe da filosofia? Acredito que não seria uma ideia bem recebida.
\begin{quote}
Kant afirmou sem rodeios que os factos observáveis do mundo físico são completa e satisfatoriamente descritos e sistematizados pela " ciência propriamente dita" ; a " filosofia propriamente dita"{} nunca nos poderá dizer nada sobre eles. A função da filosofia é erigir uma superestrutura sobre a ciência, não para fazer afirmações sobre factos físicos, mas antes para atribuir " valor"{} a certos tipos de ação humana... A investigação que procura afirmações verdadeiras para além do domínio da " ciência propriamente dita"{} e tenta obter essas afirmações a partir de uma " interpretação filosófica da ciência"{} é chamada " investigação metafísica" . O resultado deste tipo de investigação é um ramo do pensamento sistemático conhecido como " metafísica" , cuja função é fornecer-nos regras para a conduta humana, para nos mostrar " o bem"\cite{frank1952origin}.
\end{quote}
Ou seja, a filosofia conforme descrita aqui trata daquilo que não são os fatos do mundo físico. Kant então ``proíbe'' a filosofia de fazer afirmações sobre o mundo físico. Antes de concluir, acredito que é interessante acompanharmos mais um trecho:
\begin{quote}
Segundo Duhem, em suma, Platão distingue três graus de astronomia: observacional, geométrica e teológica. A conceção platónica da astronomia torna clara a natureza da divisão dentro da " filosofia"{} ou da " ciência" . A astronomia observacional e geométrica constituem, em conjunto, um instrumento que permite prever as posições observáveis dos astros, mas a " astronomia teológica"{} não é necessária nem sequer útil para esse efeito. O que ela pode fazer é dar-nos uma imagem do universo que será considerada uma réplica da sociedade humana ideal; ela pode ajudar a moldar a conduta humana\cite{frank1952origin}.
\end{quote}
E se ainda há dúvidas sobre o caráter teológico e a distinção explícita entre ciência e filosofia desde Platão, podemos consultar as próprias palavras de Platão:
\begin{quote}
Existem três graus de conhecimento. O grau mais baixo é o conhecimento pela observação sensorial. . . . O grau supremo é o conhecimento pelo intelecto puro... Entre o primeiro (mais baixo) e o terceiro (supremo) grau de conhecimento existe um tipo de raciocínio misto e híbrido que ocupa o grau intermediário (segundo). O conhecimento nascido desse raciocínio intermediário é o conhecimento geométrico. A esses três graus de conhecimento correspondem três graus de astronomia.

A percepção sensorial é responsável pela astronomia da observação. Esse tipo de astronomia persegue as curvas complicadas descritas pelas estrelas... e não pode fornecer nenhuma relação comensurável ao aritmético, nem nenhuma figura definida ao geômetra... Por meio do raciocínio geométrico, a mente produz uma astronomia capaz de figuras precisas e relações constantes. Essa “verdadeira astronomia” substitui os caminhos erráticos que a astronomia observacional atribuía às estrelas por órbitas simples e constantes; ela produz as aparências complicadas e variáveis que são conhecimento falso. ... É a aproximação do conhecimento pelo intelecto puro, que revela a terceira e suprema astronomia, a astronomia teológica. ... Na constância dos movimentos celestes, ela vê uma prova da existência de espíritos divinos que estão unidos aos corpos celestes\cite{frank1952origin}.
\end{quote}
Podemos perceber que os dois primeiros graus de conhecimento, os dois primeiros graus de astronomia são exatamente o que corresponde a ciência. O primeiro degrau corresponde ao conhecimento empírico, à experimentação. À coleta de dados e observação da natureza através dos sentidos. O segundo degrau corresponde à generalização e abstração dos dados observados, à construção de modelos gerais, na física em particular, à construção das leis matemáticas que generalizam os casos observados.

E deixamos então para a filosofia o último grau. O conhecimento superior na ideia dos filósofos. Mas novamente o que a filosofia aqui está discutindo, já não é mais o mundo físico, mas discutindo mais a questão do ideal humano, algo que visa moldar a conduta humana. Neste debate, uma definição interessante entre ciência e conhecimento técnico, é que o conhecimento técnico nos permite lidar com circunstâncias que nos deparamos antes, e o objetivo do pensamento científico é aplicar o passado a novas circunstâncias. Isto é, uma generalização do conhecimento produzido no passado para que consigamos ter uma capacidade de predição sobre novas coisas do futuro, conectando então o desenvolvimento da ciência como a evolução natural do conhecimento técnico do passado.

Evidentemente a história da humanidade é complicada e certamente a ciência não tem apenas uma origem, mas meu esforço aqui é resgatar esta conexão esquecida. Um possível caminho do conhecimento, que parte da tecnologia, passa pela ciência e culmina na ``filosofia'', considerando-a uma unificação no conhecimento pode ser tracejado. É uma possibilidade, mas isto torna a filosofia uma espécie de ciência universal, e em grande parte, retoma nossas definições anteriores da filosofia como uma concepção de mundo. Neste caso, uma concepção de mundo moldada pelo conhecimento científico.

Segundo Philipp, a primeira grande cisão entre filosofia e ciência moderna ocorre com a astronomia. O filósofo e astrônomo, até então a mesma pessoa, teria se visto obrigado a escolher uma das duas opções\cite{frank1952origin}:
\begin{itemize}
\item Manter os princípios “inteligíveis” e conformar-se à falta de concordância com os dados observados.
\item Desistir da compreensão do mundo apenas pela razão e ficar satisfeito com uma hipótese “ininteligível” que concorda com os fatos.
\end{itemize}
É evidente que a filosofia opta pelo primeiro caminho, priorizando os princípios, e a ciência pelo segundo buscando concordar com os fatos. Neste contexto, “inteligível” não significa simplesmente compreensível no sentido comum, mas aquilo que pode ser deduzido da razão pura como necessário e evidente, obedecendo a princípios racionais considerados universais e autoexplicativos. Um mundo inteligível é aquele que “faz sentido por si”, independentemente da experiência, como no ideal clássico da filosofia antiga. 

O “ininteligível” então não designa algo absurdo ou irracional, mas uma hipótese que não pode ser justificada por princípios racionais evidentes, sendo aceita apenas por sua concordância com os fatos observados. Como exemplo, um princípio inteligível no sentido clássico é a ideia de que as órbitas dos planetas deveriam ser círculos perfeitos, já que o círculo era concebido como a forma geométrica perfeita e, portanto, a mais adequada para descrever o movimento dos corpos celestes. Em contraste, a hipótese ininteligível e apoiada pelos dados, é a de que as órbitas planetárias possuem excentricidades, isto é, são elípticas. Essa hipótese concorda com os dados observacionais, mas não decorre de nenhum princípio racional autoevidente sobre perfeição, harmonia ou necessidade. A partir disso, Philipp destaca dois critérios para checar a validade de uma declaração\cite{frank1952origin}:
\begin{itemize}
\item Pode ser derivado logicamente de uma declaração auto-evidente;
\item Pode ser tratado como uma hipótese que podemos derivar consequências e testar com fatos.
\end{itemize}
A primeira consiste então em um critério “filosófico” e o segundo “científico”. Aqui busco adotar sempre que possível o critério científico.

\subsection{Superação do capitalismo}

Já comentamos sobre a filosofia marxista (\textit{vulgo} materialismo dialético), um próximo tema importante é o socialismo marxista (\textit{vulgo} socialismo científico). Socialismo científico pode ser entendido como nada mais sendo do que a investigação e a prática acerca da superação do capitalismo. Consequentemente, a construção de uma sociedade pós-capitalista (\textit{vulgo} comunista). Colocado dessa forma, o socialismo, como um dos componentes do marxismo não deve ser entendido de forma simplista como um mero modelo de estado, nem como sendo apenas uma etapa de transição. De fato, Marx nunca falou sobre socialismo neste sentido, o que há é uma distinção do comunismo em fase inferior e fase superior, assim como menções à ditadura do proletariado, mas nenhuma destas etapas recebe na obra marxiana o nome do socialismo.
\begin{quote}
Entre a sociedade capitalista e a sociedade comunista medeia o período da transformação revolucionária da primeira na segunda. A este período corresponde também um período político de transição, cujo Estado não pode ser outro senão a ditadura revolucionária do proletariado.

\ldots{}

Do que se trata aqui {[}- a fase inferior do comunismo -{]} não é de uma sociedade comunista que se desenvolveu sobre sua própria base, mas de uma que acaba de sair precisamente da sociedade capitalista e que, portanto, apresenta ainda em todos os seus aspectos, no econômico, no moral e no intelectual, o selo da velha sociedade de cujas entranhas procede\ldots Estes defeitos, porém, são inevitáveis na primeira fase da sociedade comunista

\ldots{}

Na fase superior da sociedade comunista, quando houver desaparecido a subordinação escravizadora dos indivíduos à divisão do trabalho e, com ela, o contraste entre o trabalho intelectual e o trabalho manual; \ldots{} só então será possível ultrapassar-se totalmente o estreito horizonte do direito burguês e a sociedade poderá inscrever em suas bandeiras: De cada qual, segundo sua capacidade; a cada qual, segundo suas necessidades\cite{marx1875gotha}.
\end{quote}
Já em Engels, quando se fala de socialismo científico em “Do Socialismo Utópico ao Socialismo Científico”, refere-se ao socialismo como uma ciência, no sentido amplo. Como um conjunto de conhecimentos e práticas relacionadas à investigação acerca da superação do capitalismo. Nunca como um modelo de sociedade. É neste ponto que se inserem algumas das maiores contribuições de Lênin. Sua obra aprofunda a teoria marxista ao formular uma concepção de partido revolucionário que atua como instrumento consciente da transição. Também gostaria neste assunto de destacar a contribuição do Paul Cockshott em “Towards a New Socialism” onde ousa propor um possível modelo de sociedade comunista.
\begin{quote}
Em termos de instituições democráticas e mecanismos de planejamento eficientes, devemos reconhecer que os problemas que emergiram no caso soviético refletem certas fragilidades do marxismo clássico\ldots{} Marx, Engels e Lenin foram muito mais incisivos em suas críticas ao capitalismo do que em suas teorizações positivas sobre a sociedade socialista \ldots{} Quanto aos mecanismos de planejamento, Marx e Engels apresentaram algumas sugestões interessantes, mas estas nunca foram desenvolvidas além do nível de generalidades bastante vagas\cite{cockshott1993towards}.
\end{quote}
Esta proposta recebe o nome de cibercomunismo por alguns autores, e nada mais é do que a proposta de um modelo de comunismo onde utiliza-se os mais recentes avanços científicos e computacionais no desenvolvimento e gestão desta sociedade pós-capitalista. Não ficando restrita ao uso da computação apenas para cálculos econômicos, mas transformando a tecnologia em uma ferramenta de efetiva participação popular nas decisões políticas e econômicas desta sociedade, isto é, colocando nas mãos do trabalhador o poder de decisão sobre a produção e distribuição de bens.
\begin{quote}
A primeira coisa a esclarecer sobre a natureza deste projeto é que o planejamento comunista nada mais significa do que a regulação consciente, racional e democrática da economia\ldots{} Nesse sentido, o cibercomunismo está longe de ser uma “economia algorítmica” ou um “governo por algoritmos”. Pelo contrário, trata-se de um tipo de economia onde \ldots{} os processos de avaliação e decisão individual, bem como os de deliberação democrática, permeiam toda a operação econômica: desde a escolha pessoal da profissão ou dos meios de consumo até a definição democrática dos objetivos de desenvolvimento econômico e social, passando pelo controle social do investimento, pela promoção de talentos individuais e pela implementação de novos projetos produtivos. Não se trata, portanto, de um “supercomputador” ou de uma “inteligência artificial central” que tente “adivinhar os desejos e preferências dos indivíduos”, “prever o futuro” ou “programar a vida social”. Muito pelo contrário: o que se propõe aqui é justamente ampliar o horizonte da autonomia individual e da participação democrática com o auxílio dos avanços científico-técnicos, dentro de uma estrutura de igualdade social e plenas liberdades individuais.
\end{quote}
Sobre a superação do capitalismo um tópico que não pode ser ignorado é a revolução. O que é a revolução comunista e em que condições ela se dá. Primeiro algumas definições:
\begin{quote}
\textbf{Modo de produção}: O método de produção dos bens e serviços necessários à vida (seja para saúde, alimentação, moradia ou necessidades como educação, ciência, nutrição, etc.). O Modo de Produção é a unidade das forças produtivas e das relações de produção.

\textbf{Forças de produção}: As forças produtivas são a unidade de todo o trabalho, os instrumentos de produção (ex.: edifícios e máquinas) e os sujeitos de produção (ex.: matérias-primas). 

\textbf{Relações de produção}: As relações materiais objetivas que existem em qualquer sociedade, independentemente da consciência humana, formadas entre todas as pessoas no processo de produção social, troca e distribuição de riqueza material\cite{miaGlossary}.
\end{quote}
Por revolução eu estou falando em uma mudança nas relações sociais de produção, que consequentemente conduz a uma mudança no modo de produção. Isto se dá quando as forças de produção se desenvolvem de tal modo que as atuais relações de produção passam a ser incompatíveis com as forças de produção, servindo então como uma amarra de seu potencial. Neste momento apenas uma revolução nas relações de produção é capaz de libertar as forças de produção para o seu máximo potencial através da implementação de novas relações de produção mais compatíveis com o atual nível das forças de produção, configurando então um novo modo de produção. O materialismo histórico tem sua exposição mais clara no prefácio de Para a Crítica da Economia Política. 
\begin{quote}
O resultado geral a que cheguei e que, uma vez obtido, serviu-me de guia para meus estudos, pode ser formulado, resumidamente, assim: na produção social da própria existência, os homens entram em relações determinadas, necessárias, independentes de sua vontade; essas relações de produção correspondem a um grau determinado de desenvolvimento de suas forças produtivas materiais. A totalidade dessas relações de produção constitui a estrutura econômica da sociedade, a base real sobre a qual se eleva uma superestrutura jurídica e política e à qual correspondem formas sociais determinadas de consciência. O modo de produção da vida material condiciona o processo de vida social, política e intelectual. Não é a consciência dos homens que determina o seu ser; ao contrário, é o seu ser social que determina sua consciência. Em uma certa etapa de seu desenvolvimento, as forças produtivas materiais da sociedade entram em contradição com as relações de produção existentes, ou, o que não é mais que sua expressão jurídica, com as relações de propriedade no seio das quais elas se haviam desenvolvido até então. De formas evolutivas das forças produtivas que eram, essas relações convertem-se em entraves. Abre-se, então, uma época de revolução social. A transformação que se produziu na base econômica transforma mais ou menos lenta ou rapidamente toda a colossal superestrutura... Do mesmo modo que não se julga o indivíduo pela ideia que de si mesmo faz, tampouco se pode julgar uma tal época de transformações pela consciência que ela tem de si mesma. é preciso, ao contrário, explicar essa consciência pelas contradições da vida material, pelo conflito que existe entre as forças produtivas sociais e as relações de produção. 

\cite{marx2008critica}
\end{quote}
Em segundo lugar, eu acredito que a gente vive uma situação particularmente diferente de todos países que que sofreram revolução. Próximo a época de revolução, estes países viveram uma situação de intensa instabilidade, e não por conta dos comunistas. Ou seja, a instabilidade que estes países se encontravam não foi produto dos comunistas, o que os comunistas fazem é guiar esta revolução. Por exemplo, a China estava ocupada militarmente pelo Japão, já Cuba vivia sob a ditadura de Fulgêncio Batista, etc... Entendo que isto cria verdadeiramente uma sensação de que “não há nada a perder”. 

De certa forma, talvez a Rússia seja o país mais próximo do Brasil, já que ela não era colonizada. Mas ainda sim, o império russo estava envolvido em uma Guerra Mundial e sequer encontrava-se ainda com um modo de produção capitalista. Ou seja, não vivia sob um “capitalismo democrático” como o Brasil. Mesmo se tentássemos ignorar os efeitos da Primeira Guerra Mundial, devemos lembrar que a assembleia constituinte formada entre a revolução de fevereiro (revolução burguesa) e a revolução de outubro (revolução socialista) constava com mais de 20\% dos deputados eleitos pertencendo ao partido bolchevique\footnote{Das 767 cadeiras, os bolcheviques eram o segundo maior partido, ocupando 183.}. Isto é, em um primeiro momento houve uma revolução burguesa, e no segundo momento, anterior à revolução socialista, os comunistas gozavam de cada vez mais popularidade. No momento não temos nenhum dos dois cenários no Brasil.

Hoje o Brasil é relativamente estável politicamente. Além disso, o capitalismo aprendeu a ceder algumas migalhas em prol da estabilidade\cite{piketty2014capital}. Como resultado, a classe média\footnote{Classe baixa, média e alta não são sinônimos de proletariado, pequena burguesia e burguesia, mas há entre elas uma correspondência aproximada. Sobre o desenvolvimento da classe média no século passado, “O Capital no Século XXI” de Thomas Piketty é uma interessante leitura.} cresceu e se tornou uma das principais forças políticas, chegando a rivalizar com a classe baixa em termos de tamanho (a depender da métrica utilizada) e possuindo mais poder político. Se a pequena burguesia já era uma classe vacilante na época de Marx, ela possivelmente é ainda mais nos dias de hoje devido ao poder econômico e político que ganhou nas últimas décadas. 
\begin{quote}
A ascensão do capitalismo a partir do longo século XVI está associada a uma queda nos salários para níveis abaixo da subsistência, uma deterioração na estatura humana e um aumento na mortalidade prematura. \ldots{} Onde houve progresso, melhorias significativas no bem-estar humano começaram apenas por volta do século XX. Esses avanços coincidem com a ascensão de movimentos políticos anticoloniais e socialistas\cite{sullivan2023capitalism}.
\end{quote}
Então como eu vejo a revolução sendo possível? O caminho que consigo visualizar é através da proposta de levar aos limites a democracia burguesa. Eu não vejo a possibilidade de convencer a população a expropriar os capitalistas enquanto as pessoas não defenderem sequer a existência de um salário mínimo, eu não vejo como vamos convencer a planificar a economia se as pessoas não defendem nem mesmo a existência de um imposto progressivo, então eu acredito que ainda existe um grande trabalho a ser feito dentro dos limites da democracia burguesa.

Como lidar com uma população que é fortemente povoada por uma classe média com poder político e vacilante? Eu defendo que devemos seguir com certa precaução o exemplo do Salvador Allende, usar as vias legais tanto quanto possível. Se pudermos avançar em reformas, não há porque nos negarmos. Isso serve de propaganda das ideias comunistas de duas formas. Primeiro, quando os projetos são aprovados e executados serve como uma demonstração que nosso projeto de sociedade é de fato superior. E segundo, quando reprovado serve como uma de denúncia sobre a perseguição que as ideias comunistas sofrem, mesmo que sejam do interesse da maioria dos trabalhadores. 
\begin{quote}
Sob as condições de que a classe dominante em um país capitalista esteja disposta a permitir que o Partido Comunista opere legalmente, o referido partido não rejeita a oportunidade. Afinal, o principal interesse de tal partido é elevar o nível de consciência do proletariado e de outras pessoas e organizá-las. Reformas também podem ser conquistadas de tempos em tempos\cite{sison2025basic}.
\end{quote}
É importante ressaltar que não tenho por meta uma nem uma revolução armada e nem um cenário de instabilidade. Não tenho por objetivo a violência de nenhuma forma. No melhor dos cenários o comunismo é atingido via reformas legais. Minha principal divergência com partidos sociais democratas se dá principalmente pelo horizonte reformista apresentado no discurso. Neste processo, eu defendo que o partido comunista tenha abertamente a intenção de superação do capitalismo, isto é, um horizonte revolucionário. As propostas devem seguir esse caminho, mesmo quando sabemos que serão rejeitadas, a rejeição deve servir de propaganda. O discurso deve se manter transparente, o objetivo não é tornar o capitalismo mais dócil, mas superá-lo. Ainda que, enquanto não superado, seja útil promover também projetos reformistas que garantam um alívio na classe trabalhadora.

Através de uma postura propositiva que inclua projetos de políticas públicas realistas, a atuação política das organizações comunistas devem permitir que as pessoas vislumbrem uma nova possibilidade de vivência e percebam que é possível ir além do capitalismo. De fato, nada disso é totalmente novo. Apenas desejo ressaltar a importância de utilizarmos os meios legais até onde for possível para promover o desenvolvimento da qualidade de vida da classe e das forças de produção, até atingirmos um ponto que uma mudança nas relações de produção torna-se uma possibilidade concreta. 

Isso também implica que a questão central é a disputa da consciência de classe. Hoje em dia, as forças militares do estado são mais fortes do que em qualquer outro momento da história. O poder bélico estatal não tem adversários à altura que não seja outro poder militar estatal. Mais do que nunca, é necessário o apoio da população, é a população que deve fazer resistência contra o ataque da burguesia. População, do qual, inclusive as forças policiais e militares fazem parte.

Tomando como exemplo a Revolução de Outubro: ``os rebeldes Bolcheviques encontraram pouca resistência e foram capazes de penetrar no edifício sem muita dificuldade e tomá-lo. Na sua maioria, a revolta em Petrogrado ocorreu sem derramamento de sangue e numa atmosfera de normalidade geral na capital. Estima-se que apenas cinco marinheiros e um soldado entre os atacantes foram mortos no assalto e que os defensores não sofreram baixas fatais.\cite{wikipedia2026revolucao}''

Esta necessidade de que a revolução seja feita pela massa de trabalhadores deve estar sempre muito clara. Quem faz a revolução é, e sempre foi a classe trabalhadora. O partido comunista é apenas o guia da revolução. Mas não é o partido que faz a revolução. Para isso o partido deve estar inserido na classe trabalhadora deve se tornar capaz de reconhecer seus objetivos e anseios, e assim conseguir apresentar propostas efetivas, que são tomadas pela classe trabalhadora como um guia.
\begin{quote}
Para triunfar, a revolução deve ser dirigida pela classe operária \ldots{} A moralidade revolucionária consiste em unir-se às massas, em ter confiança nelas, em compreendê-las, em escutar atentamente suas opiniões. Por suas palavras e seus atos, os membros do Partido, da União da Juventude trabalhadora e os quadros ganham a confiança, a admiração e o afeto do povo, realizam a união estreita das massas em torno do Partido, organizam-nas, fazem a propaganda entre elas, mobilizam-nas, para aplicar com entusiasmo a política e as resoluções do Partido \cite{hochiminh1958moralidade}.
\end{quote}
Os trabalhadores devem desenvolver uma consciência acerca do poder que a própria classe possui, uma consciência acerca das contradições do capitalismo, e assim devem se tornar conscientes da necessidade de superação do capitalismo. Apenas através desse processo de educação que entendo que a classe trabalhadora seja capaz de aderir ao projeto político comunista e termine então por defender a revolução comunista.

Quando se discute sobre a via legal como uma ferramenta para a construção do comunismo, torna-se impossível não discutir o significado da eleição de candidatos que se declaram socialistas ou comunistas. Pensando nos acontecimentos recentes, podemos refletir sobre a eleição nos EUA, especificamente para prefeito de Nova York, onde algumas coisas me chamaram a atenção. Eu quero deixar claro que não o considero revolucionário, comunista, nem sequer socialista. Grande parte da campanha de Zohran Mamdani foi baseada em diferenciar o que ele entende de socialismo do que ele entende de comunismo, deixando claro que não é comunista e limitando sua defesa do socialismo a uma redução da desigualdade sem superação do modo de produção capitalista. Ao mesmo tempo que se vende como uma campanha socialista, também promove uma campanha que se apoia fortemente no anticomunismo para traçar essa linha de separação entre ambos. Uma consequência desta prática, é que apaga-se a separação do capitalismo e socialismo, alterando o próprio sentido de socialismo. Isto faz com que a luta-anticapitalista sob esta bandeira do socialismo deixe de ser sobre a socialização dos meios de produção, e se torne meramente um capitalismo com menor desigualdade através de algumas políticas de redistribuição de renda.
\begin{quote}
Em vez do lema conservador de: “Um salário justo por uma jornada de trabalho justa!”, deverá inscrever na sua bandeira esta divisa revolucionária: “Abolição do sistema de trabalho assalariado!”\cite{marx1865salario}.
\end{quote}
Mas dito isso, eu ainda acredito que há o que aprender desta experiência. Afinal é a maior cidade do mundo, e me parece evidente que a principal preferência burguesa nestas eleições saiu derrotada. De forma alguma eu estou negando o poder e influência da ideologia dominante, até porque quem venceu não foge totalmente dela, o ``socialismo'' venceu dentro dos limites da defesa do capitalismo.

Mas ainda sim, demonstra como a política é complexa, e sob certas condições, mesmo uma campanha com menos dinheiro pode derrotar o principal candidato da burguesia. Evidentemente a disponibilidade de recursos financeiros e a postura adotada do candidato diante da classe burguesa ainda é muito importante e claramente apontam tendências sobre quem serão os favoritos, mas sobra espaço para um certo grau de manobra em determinados contextos.

Eu entendo que estas eleições recentes não foram uma derrota da burguesia pelo proletariado, o candidato eleito não pode trazer grandes mudanças estruturais pois independente das ideias que pessoalmente traz, sua atuação está limitada pelo próprio funcionamento da estrutura capitalista. Mas eu vejo como uma derrota da grande burguesia para a pequena burguesia, e isto é significativo. Evidentemente que a burguesia pode impedi-lo de implementar todas as propostas fazendo com que o episódio se converta em uma derrota para a pequena-burguesia (e possivelmente para toda classe trabalhadora), mas isso não apaga o que foi feito. Como diz Antonio Candido: ``O que se pensa que é a face humana do capitalismo é o que o socialismo arrancou dele\cite{tavares2021socialismo}''.

\subsection{Guerra cultural}

Apesar da ênfase da minha pesquisa atual ser sob a esfera econômica da sociedade, a esfera cultural não pode ser negligenciada. O trabalho é, para Marx, uma atividade teleológica, e de certa forma, este é o grande diferencial do trabalho humano em relação às atividades de outros animais. No primeiro volume de O Capital, Marx distingue o trabalho do tecelão à atividade de uma aranha escrevendo:
\begin{quote}
Nós supomos o trabalho numa forma em que ele pertence exclusivamente ao homem. Uma aranha realiza operações que se assemelham às do tecelão e uma abelha, através da construção dos seus alvéolos de cera, envergonha muitos mestres-de-obras humanos. O que, porém, de antemão distingue o pior mestre-de-obras da melhor abelha é que ele construiu o alvéolo na sua cabeça antes de o construir em cera. No fim do processo de trabalho obtém-se um resultado que, no começo do mesmo, já na ideia do operário, portanto, já idealmente, se achava presente\cite{marx2013capital1}.
\end{quote}
Ou seja, um dos principais aspectos do trabalho humano é que quando realizamos um trabalho, nós o realizamos seguindo um plano que elaboramos em nossa mente. Nós agimos baseados neste plano visando atingir um fim específico. Inevitavelmente o fim que queremos atingir vai determinar qual é o plano que vamos adotar, e o plano que adotamos vai determinar como vamos agir. Não precisamos aqui entender determinar’ de forma muito restrita, devemos ler no sentido de ``limitar'', ou seja, os fins que queremos atingir limitam que estratégias podemos desenvolver e consequentemente também como podemos agir, ainda que mais de uma estratégia ou ação esteja disponível.

Ideia semelhante é abordada sob uma perspectiva antropológica em um dos melhores livros que eu li, o livro de Agustín Fuentes chamado “Why We Believe: Evolution and the Human Way of Being”\cite{fuentes2019why}. Neste livro Agustin não apenas reconhece esta capacidade tão distintivamente humana de trabalhar, mas também associa esta forma de proceder com a nossa própria capacidade de imaginar no seu sentido mais amplo possível.

A partir desta abordagem antropológica, Agustin demonstra que nossa capacidade de imaginar é uma habilidade vital para a sobrevivência de nossa espécie, certamente resultado do próprio processo de seleção natural. A partir desta leitura somos expostos à tese de que a capacidade de imaginar coisas possíveis e coisas impossíveis é uma só. Poderíamos estender a reflexão ao fato de que, na prática, muitas coisas que imaginamos ser possível, só estaríamos em condições de descobrir que estávamos errados após tentarmos realizar esta imaginação na vida real. Ou seja, é a prática real que distingue o que pensamos entre o que é possível e o que não, mas o processo de imaginar é o mesmo.

Esta capacidade então de imaginar e acreditar que nasce com a própria atividade do trabalho se estende então para todas esferas da vida humana. É fácil perceber como ela é a mesma capacidade humana que nos permite criar ficção e a arte como um todo. O gênio criativo humano que nos permite desenvolver teorias científicas elaboradas e tecnologias modernas, que permite avançarmos na compreensão da natureza é fundamentalmente o mesmo que nos permite produzir arte.

A esta altura, você pode estar se perguntando aonde eu quero chegar. Em “O que fazer?” Lenin escreveu: “É preciso sonhar!”\cite{lenin1902quefazer}. A ideia que eu quero formular é relativamente simples: para que o ser humano execute seu trabalho de forma adequada, ele precisa ter um plano de como realizar o trabalho para atingir um determinado fim, e para isso ele precisa exercer seu poder de imaginação. Não é possível executar um trabalho de forma adequada sem conseguir antecipar antes mentalmente qual é o resultado deste o trabalho.

O impacto disto é profundamente importante e acredito, que muitas vezes ainda subestimado. A construção de uma nova sociedade pós-capitalista, o objetivo final de todo mundo movimento comunista, só pode ser realizado como fruto do trabalho humano. E sendo assim, como todo trabalho humano, este é um esforço que acontece de forma orientada, buscando concretizar uma sociedade que foi previamente planejada, ainda que o planejamento não tenha sido elaborado de forma detalhada, a sociedade futura foi previamente imaginada. Não é surpreendente então a afirmação de que para construir uma nova sociedade, é preciso primeiro que sejamos capazes de imaginar uma nova sociedade.

É evidente que como este é um movimento que inclui milhões de seres humanos anti-capitalistas buscando tornar real sua própria ideia de uma nova sociedade, em um confronto ideológico com outros tantos milhões de defensores do capitalismo que que buscam concretizar também sua própria ideia de sociedade (e ainda que estas pessoas defendem a manutenção do capitalismo, eles precisam projetar esta sociedade no futuro através da imaginação e trabalhar para que o capitalismo do futuro corresponda aquela ideia que ele tem de como o capitalismo deveria ser neste futuro), o resultado final deste confronto de intenções não é simples. O resultado final não será exatamente a projeção mental de nenhum indivíduo isoladamente, como Engels já bem colocou em uma carta a Joseph Bloch de 21 de setembro de 1890:
\begin{quote}
a história faz-se de tal modo que o resultado final provém sempre de conflitos de muitas vontades individuais, em que cada uma delas, por sua vez, é feita aquilo que é por um conjunto de condições de vida particulares; há, portanto, inúmeras forças que se entrecruzam, um número infinito de paralelogramas de forças, de que provém uma resultante --- o resultado histórico ---, que pode ele próprio, por sua vez, ser encarado como o produto de um poder que, como todo, actua sem consciência e sem vontade. Pois, aquilo que cada indivíduo quer é impedido por aquele outro e aquilo que daí sai é algo que ninguém quis. Assim, a história até aqui decorreu à maneira de um processo natural e está também essencialmente submetida às mesmas leis de movimento. Mas, de que as vontades individuais --- em que cada um quer aquilo a que o impele a sua constituição física e circunstâncias exteriores, em última instância económicas (quer as suas próprias pessoais quer as gerais-sociais) --- não alcançam aquilo que querem, mas se fundem numa média total, numa resultante comum, daí não deve, contudo, concluir-se que elas são de pôr como igual a zero. Pelo contrário, cada uma contribui para a resultante e está, nessa medida, compreendida nela\cite{engels1890bloch}.
\end{quote}
Aqui eu quero fazer um breve comentário para defender Engels de acusações de positivismo. É preciso ter em mente que esta analogia com as leis da natureza é até certo ponto apenas isto, uma analogia, Engels não está submetendo o desenvolvimento da sociedade humana à segunda lei de Newton, apenas operando uma analogia onde visa demonstrar que nossa sociedade não se encontra fora da natureza. Assim como Marx apontou na Ideologia Alemã: 
\begin{quote}
Conhecemos uma única ciência, a ciência da história. A história pode ser examinada de dois lados, dividida em história da natureza e história dos homens. Os dois lados não podem, no entanto, ser separados; enquanto existirem homens, história da natureza e história dos homens se condicionarão reciprocamente. A história da natureza, a assim chamada ciência natural, não nos diz respeito aqui; mas, quanto à história dos homens, será preciso examiná-la, pois quase toda a ideologia se reduz ou a uma concepção distorcida dessa história ou a uma abstração total dela\cite{marx2007ideologia}.
\end{quote}
Pode-se dizer inclusive que Engels adianta o desenvolvimento da econofísica conforme se daria apenas próximo dos anos 2000 onde é demonstrado matematicamente que podemos observar comportamentos análogos quando tratamos a sociedade como um sistema complexo e outros tipos de sistemas complexos tipicamente encontrados em ciências da natureza. Há um fundo científico na analogia. Mas retomando a discussão que nos propomos inicialmente, a questão central é que mesmo que o futuro da sociedade não esteja perfeitamente na mente de nenhum dos membros que a constitue, o seu desenvolvimento é resultado de um confronto de vontades individuais, que agem conforme imaginam e projetam como a sociedade deveria ser no futuro. E é com esta perspectiva que a capacidade de imaginar como a sociedade poderia ser é determinante pro desenvolvimento de como a sociedade realmente será.

Neste contexto, defendo então duas ideias: 
\begin{enumerate}
\item A construção da sociedade é fruto de diversos processos teleológicos individuais, na mesma medida que o trabalho humano é em si próprio uma atividade teleológica; 
\item A disputa cultural pelo imaginário popular é uma disputa fundamental para a construção de uma sociedade melhor.
\end{enumerate}
Este seja um bom momento de compartilhar algumas anotações do livro “Revolução, arte e cultura” de Anatoli Lunatcharski onde ele ressalta o impacto da arte na construção da sociedade. É importante também destacar que o que fazemos é reconhecer a influência da arte e da cultura na construção da sociedade em qualquer época, e não reduzir a arte do futuro a propaganda política.
\begin{quote}
As classes dominantes a utilizaram para erigir a sociedade de acordo com os seus interesses. A arte também foi utilizada como arma na luta pelas classes que foram levadas, por seu desenvolvimento histórico, a se contrapor àquelas classes dominantes. A chamada “arte pela arte” é uma arte que foge da vida, que se mantém extremamente distante dos problemas reais da vida, e até mesmo pronunciadamente os despreza; é uma arte que, ativa ou passivamente, consciente ou inconscientemente, se distancia das forças sociais; mas também é, apesar de tudo, uma força social, que às vezes serve de forma inequívoca e determinada a interesses específicos\cite{lunatcharski2022revolucao}.
\end{quote}
Mas dentro do contexto político que nos inserimos, onde visamos construir uma sociedade melhor, a arte então coerente com este projeto construir uma contra-cultura e assume um papel central:
\begin{quote}
Nossa arte não pode ser outra coisa senão uma força que exerce violenta influência no transcorrer geral da luta e da construção\ldots{} {[}consideramos{]} a arte, absoluta e completamente, como uma divisão do nosso exército de luta e de construção do socialismo.

A arte é uma grande força em nossas mãos para a realização dessa obra. A arte é capaz não só de orientar, mas também de formar. Para o artista, não se trata apenas de mostrar a toda sua classe como o mundo é hoje, mas de ajudá-la a encontrar o seu caminho na existência, de educar o homem novo. Com isso pode acelerar o ritmo de desenvolvimento da realidade, e ele pode construir, através da criação artística, um centro ideológico que se encontre acima dessa realidade, que a eleve, que permita uma perspectiva do futuro e, com isso, acelerar o ritmo\cite{lunatcharski2022revolucao}.
\end{quote}
Grande parte das minhas ideias sobre cultura e arte são inspiradas no trabalho de Bakhtin. Sendo assim acho pertinente uma pequena exposição sobre o que entendo das ideias deste pensador antes de avançarmos, tomando então a mediação da vida social através dos signos como ponto de partida da nossa discussão. Bakhtin aprofunda a discussão marxista quanto à relação entre matéria e consciência, a partir do pressuposto de que a ideologia é um fenômeno social.
\begin{quote}
Todo signo é ideológico; a ideologia é um reflexo das estruturas sociais; assim, toda modificação da ideologia encadeia uma modificação da língua\cite{bakhtin1929marxismo}.
\end{quote}
Isso tendo em mente o que ideologia significa para Bakhtin:
\begin{quote}
Para o Círculo de Bakhtin, a palavra ideologia tem significado diferente daquele usado por parte da tradição marxista\ldots{} a palavra ideologia é usada em geral para designar o universo dos produtos do “espírito humano”, aquilo que algumas vezes é chamado por outros autores de cultura imaterial ou produção espiritual\ldots{} Assim, ideologia para o Círculo de Bakhtin, abrange um grande universo: a arte, a filosofia, a ciência, religião, ética, política. Todo produto ideológico parte de uma realidade (natural ou social), possui um significado e remete a algo que lhe é exterior, ou seja, é um signo\cite{david2017ideologia}.
\end{quote}
Nesta concepção a ideologia pode ser encontrada nos signos que utilizamos nesta mesma vida social e existe de forma exterior a nós, como um produto de nossa atividade social. Estes signos por sua vez tem a importante característica de que estes signos estão sempre inseridos em uma luta ideológica. Dessa forma, o primeiro ponto importante é que devemos entender que nossa própria consciência emerge da nossa vida social mediada por signos. 

O segundo ponto importante é que devemos lembrar que o signo não tem uma existência meramente abstrata, ele precisa ter uma existência material, ele precisa existir como algo que possa ser percebido e compartilhado, mesmo que seja um som ou um registro escrito. Dessa forma a classe dominante ao controlar os meios de produção materiais, também controla a produção e compartilhamento de signos, e de maneira geral, a própria cultura hegemônica.

Um terceiro ponto, é que ao considerar que todo signo estar inserido em uma disputa ideológica, isto é, quando selecionamos um signo realizamos também um julgamento de valor, então estamos implicando no fato de que forma e conteúdo estão relacionados de maneira indissociável, já que uma vez que trocamos a forma, trocamos o signo, e consequentemente, pode ser estejamos expressando um diferente valor. Para ficar claro, por signo aqui podemos entender como o produto da unidade entre pelo menos um significante, um significado e um valor. Isto é, a palavra “cachorro” é um signo, o significado é o próprio animal que nos referimos, o símbolo gráfico a partir do qual lemos a palavra é o significante se for escrito, e se for falado, então é o som correspondente à palavra que é o significante. Além disso a escolha por um signo carrega em si também um julgamento de valor quanto ao que 

Podemos encontrar diversas disputas políticas quanto ao uso de diferentes palavras para se referir a mesma coisa, principalmente no contexto das lutas sociais promovidas pela população historicamente marginalizada, como por exemplo, a escolha entre as palavras “índio” e “índigena”, ou “tribo” e “nação”. Em uma conversa oral, até mesmo uma mudança na entonação em uma mesma palavra pode sugerir diferentes interpretações para a mensagem final. 

É importante ter em mente que a construção do significado, valor e do próprio signo em si se dá através da vida social coletiva de uma comunidade. O que implica também que diferentes signos podem ser percebidos de forma diferente em diferentes comunidades. Não há atividade social neutra. A propagação ou combate de determinadas ideologias se faz presente em todos aspectos da vida social.
\begin{quote}
\ldots{} a “ideologia do cotidiano”, que se exprime na vida corrente, é o cadinho onde se formam e se renovam as ideologias constituídas\ldots{} A ideologia do cotidiano constitui o domínio da palavra interior e exterior desordenada e não fixada num sistema, que acompanha cada um dos nossos atos ou gestos e cada um dos nossos estados de consciência. Considerando a natureza sociológica da estrutura da expressão e da atividade mental, podemos dizer que a ideologia do cotidiano corresponde, no essencial, àquilo que se designa, na literatura marxista, sob o nome de “psicologia social”\cite{bakhtin1929marxismo}.
\end{quote}
Dentre os signos existentes, a palavra ocupa uma posição especial, pois ela é um instrumento da consciência e é devido a esse papel que a “\ldots{} palavra funciona como elemento essencial que acompanha toda criação ideológica, seja ela qual for.” Isto se dá pelo fato de que a consciência tem o poder de abordar cada signo, independente da sua natureza, de modo verbalmente. Então a palavra se torna “\ldots o objeto fundamental do estudo das ideologias”. Por isso então:
\begin{quote}
A única maneira de fazer com que o método sociológico marxista dê conta de todas as profundidades e de todas as sutilezas das estruturas ideológicas “imanentes” consiste em partir da filosofia da linguagem concebida como filosofia do signo ideológico\cite{bakhtin1929marxismo}.
\end{quote}
Um grande progresso que se faz a partir de Bakhtin é diluir a fronteira entre a atividade mental como atividade interior e a comunicação como atividade exterior. Ambos são tratados como uma mesma atividade, onde a atividade mental interior é vista como uma “conversa interna”. A diferença entre ambos se dá mais pelo grau que pela qualidade, a atividade mental tende a ser mais confusa e fragmentada, ela ganha clareza à medida que passa a se transformar em uma comunicação. Mesmo quando organizamos nossas ideias como uma auto explicação, ela já se aproxima desta dimensão de exteriorização da atividade mental mais organizada. 

Essa forma de entender a consciência forma uma base radical de pensar na vida social humana. Se os signos utilizados na comunicação tem uma origem social. Isso imprime uma natureza social na nossa própria atividade mental interior. Ficamos então impossibilitados de que estejamos, a qualquer momento, totalmente isolados das influências culturais da sociedade. E por sua vez, a natureza social dos signos também os vincula a uma realidade material a partir da qual essa realidade social é construída. Assim sendo, nossa própria atividade mental interior encontra-se em uma relação direta com a realidade material no qual estamos inseridos e que nos fornece os signos que internalizamos.

Nossa existência então, como animais sociais, se dá tanto em uma dimensão material onde a vida material se desenrola, quanto em uma dimensão simbólica (ou dimensão ideológica) onde a vida social acontece. Ambas dimensões relacionam-se entre si de forma dialética, isto é, ambas dimensões afetam e são afetadas por sua contraparte.

Resumindo:
\begin{itemize}
\item Não existe pensamento sem signos;
\item Não existe signo que não seja social;
\item Logo, a consciência individual é sempre atravessada pela vida social.
\end{itemize}
Podemos nos afastar então tanto de uma leitura economicista da sociedade (ou seja, ver a cultura como um mero reflexo do capitalismo) quanto de uma leitura culturalista (ver a cultura como uma dimensão isolada das bases materiais da sociedade). Mas não podemos deixar de perceber que a cultura cumpre um papel vital na sustentação do capitalismo.

Caminhando para a conclusão dessa sessão, estamos então em uma situação que é interessante definir o que queremos dizer com a palavra cultura: um conjunto historicamente produzido de signos, valores\footnote{Sobre valores capitalistas e anti-capitalistas uma leitura interessante é ``Como ser anticapitalista no século XXI?'' de Erik Olin Wright\cite{wright2019anticapitalista}.} e práticas que mediam toda atividade social, organizando a experiência humana em torno de uma concepção de mundo. Dito isso, é útil relembrarmos o que é concepção de mundo:
\begin{quote}
Concepção de mundo: abrange a totalidade das opiniões e concepções de um indivíduo ou sociedade sobre o mundo, sobre nós mesmos como seres humanos e sobre a vida e a posição dos seres humanos no mundo\cite{worldview2023}.
\end{quote}
A proposta revolucionária do marxismo consiste no seu objetivo em promover a superação do modo de produção capitalista. Uma vez que o que define se o modo de produção como capitalista são as relações de produção, a revolução socialista não é outra coisa senão uma mudança nas relações de produção. Assim sendo, as relações de produção ocupam uma posição central no debate. Uma vez entendido isso, para avançar na reflexão podemos sintetizar três ideias já expostas:
\begin{itemize}
\item De acordo com a própria definição das relações de produção, os seres humanos não podem produzir fora de uma estrutura social.
\item De acordo com Bakhtin toda atividade social é mediada por signos que nunca são neutros.
\item De acordo com nossa definição, a cultura é um conjunto de signos historicamente produzidos.
\end{itemize}
Com isso em mente, podemos perceber como a cultura é parte da estrutura social no qual as relações de produção - e consequentemente o modo de produção - existem. Ainda sim, vale enfatizar, a cultura não deve ser vista como autônoma, ela é determinada pelo modo de produção, porém, por “determinada”, o que queremos dizer é que o modo de produção impõe um limite nas possibilidades de desenvolvimento cultural que podem emergir nesta sociedade, não que não existam possibilidades.

Nem toda relação social é uma relação de produção. Relações de produção possuem um significado mais estrito, podemos ler na definição dada pelo glossário da MIA que “a base das relações de produção é a propriedade dos meios de produção”\cite{miaGlossary}. Ou seja, o modo de produção determina (limita) a cultura, porém também depende da existência de uma cultura hegemônica correspondente coerente com a sua própria existência. Não há um modo de produção que exista a parte de qualquer cultura.

Como toda atividade social é mediada por signos que nunca são neutros, e estes são produzidos socialmente dentro de uma dada cultura hegemônica, a cultura pode ser entendida como um dos mecanismos responsáveis por manter o modo de produção capitalista estável de forma “democrática”. É impossível superar o capitalismo sem mexer na sua base econômica, mas sua estabilidade também demanda uma cultura propícia para tal.

E aqui consiste uma das principais características do capitalismo, o senso de liberdade e escolha democrática construído pela cultura hegemônica. O fato é que grande parte dos trabalhadores aceitam se submeter à exploração capitalista “voluntariamente”, pois esta exploração foi (através da cultura hegemônica) naturalizada.

A opressão em outros modos de produção pré-capitalistas se dava através de mecanismos extra-econômicos, ou seja, a opressão se dava muitas vezes através da força bruta. O nascimento do capitalismo é marcado pela separação das esferas político-jurídica e econômica, de modo que há também uma modificação nos meios empregados para a exploração, de meios extra-econômicos para meios econômicos. Ellen Wood discute brilhantemente esta etapa de transição em seu livro ``A Origem do Capitalismo''\cite{wood2001origem}. Esta distinção entre os diferentes meios de coerção adotado pelas sociedades pré-capitalistas e capitalistas está presente desde Marx:
\begin{quote}
No evolver da produção capitalista desenvolve-se uma classe de trabalhadores que, por educação, tradição e hábito, reconhece as exigências desse modo de produção como leis naturais e evidentes por si mesmas \ldots{} a coerção muda exercida pelas relações econômicas sela o domínio do capitalista sobre o trabalhador. A violência extraeconômica, direta, continua, é claro, a ser empregada, mas apenas excepcionalmente\cite{marx2013capital1}.
\end{quote}
Podemos comparar esta situação imposta pelo capitalismo, com uma situação híbrida, onde o capitalismo ainda não se tornou o modo de produção dominante.
\begin{quote}
Basta, aqui, uma simples alusão a formas híbridas, em que o mais-valor não se extrai do produtor por coerção direta e que tampouco apresentam a subordinação formal do produtor ao capital. Nesses casos, o capital ainda não se apoderou diretamente do processo de trabalho\cite{marx2013capital1}.
\end{quote}
Fica então evidente o contraste entre a coerção direta de formas pré-capitalistas e a coerção indireta no capitalismo. Merece uma atenção especial o fato de que a classe de trabalhadores aceita o modo de produção através da “educação, tradição e hábito”:
\begin{quote}
Não basta que as condições de trabalho cristalizem em um dos polos como capital e no outro polo contrário como homens que não têm mais nada que vender a não ser sua força de trabalho. Não basta tampouco obrigar estes a se vender voluntariamente. No transcurso da produção capitalista, vai sendo formada uma classe operária que, por força de educação, de tradição, de costume, se submete às exigências deste regime de produção como às mais lógicas leis naturais. A organização do processo capitalista de produção já desenvolvido vence todas as resistências, a existência constante de uma superpopulação relativa mantém a lei da oferta e da demanda de trabalho em concordância com as necessidades de exploração do capital, e a pressão surda das condições Econômicas sela o poder de mando do capitalista sobre o operário. Ainda é empregada, de vez em quando, a violência direta, extraeconômica; mas só em casos excepcionais. Dentro do transcurso natural das coisas, já pode deixar-se o operário a mercê das 'leis naturais da produção', isto é, entregue ao predomínio do capital, predomínio que as próprias condições de produção engendram, garantem e perpetuam\cite{marx2013capital1}.
\end{quote}
Esta ideia reaparece no mesmo livro de Ellen Wood, quando é feito um comentário sobre as contribuições de Thompson.
\begin{quote}
Em sua obra o desenvolvimento do capitalismo ganha vida não apenas como um processo de proletarização, particularmente em seu clássico A formação da classe trabalhadora inglesa (1963), mas também como um confronto vivo entre os princípios do mercado, as práticas e valores alternativos\cite{wood2001origem}.
\end{quote}
Ou seja, para que o capitalismo possa surgir, temos um confronto vivo entre diferentes valores e práticas. Se lembrarmos que definimos cultura como um conjunto de signos, valores e práticas, podemos presumir que foi necessário a formação de uma nova cultura hegemônica. Para complementar esta discussão, também é possível encontrar um exemplo quanto à questão dos signos no mesmo livro. Enquanto discorre sobre o desenvolvimento do capitalismo Wood chama a atenção do leitor para o fato de que a palavra “melhorar” em sua acepção original significava literalmente “fazer alguma coisa com vistas ao lucro monetário”, então propõe a seguinte reflexão:
\begin{quote}
\ldots{} essa palavra foi adquirindo um significado mais geral, no sentido como a conhecemos hoje (e seria interessante pensar nas implicações de uma cultura em que a palavra correspondente a “tornar melhor” enraíza-se no termo que corresponde a lucro monetário)\cite{wood2001origem}.
\end{quote}
Na mesma linha de raciocínio, há também uma breve discussão análoga sobre as mudanças no significado da palavra “produtor” com o advento do capitalismo. Assim, não deve nos surpreender que Thompson, segundo Wood novamente, teria dito que um dos momentos axiais da formação do capitalismo seja o “o processo em que foram forjados um novo proletariado e uma nova cultura da classe trabalhadora.” 

A exploração no capitalismo nasce então da atuação das forças econômicas, da necessidade que o trabalhador tem de vender a sua força de trabalho. Ainda que em teoria, qualquer um pode se tornar explorador, é preciso apenas dinheiro, e não mais depende de algo imutável como descendência real, é principalmente a estrutura social que mantém essa divisão entre explorados e exploradores.

Isto resulta na constatação de que um dos grandes focos de poder no capitalismo, se dá na dimensão da cultura. Isto é, na construção dos nossos valores e gostos, na construção do nosso imaginário. Uma das grandes ferramentas de controle e organização do capitalismo é sua capacidade de criar uma cultura que perpetua e justifica a exploração. A força militar surge de forma secundária principalmente em momentos de crise, mas não é a principal força atuante no cotidiano para a maior parte da população, ainda que inegavelmente se faça presente.

Por isso é preciso ter em mente que a indústria no capitalismo não se reduz a “indústria material”, aquela que tipicamente pensamos quando pensamos em indústria, ou seja, aquela indústria que produz bens materiais. Mas a indústria capitalista inclui, e não com menor importância, a indústria cultural, isto é, a indústria que produz cultura. Como exemplo desta indústria, podemos lembrar da indústria de Hollywood estado-unidense exportando uma cultura atrelada ao “American Way of Life”, ou então o fenômeno global de popularização da cultura sul-coreana chamada de onda coreana “Hallyu”, também poderíamos citar o Japão e o soft-power construído a partir da indústria de animes e mangás. 

Podemos trabalhar então com uma definição de cultura hegemônica:\textbf{ }refere-se ao domínio de um grupo cultural sobre outros dentro de uma sociedade. Esse domínio é estabelecido e mantido por meio de instituições culturais, como a mídia, a educação e a religião. A hegemonia cultural molda percepções, normas e valores, muitas vezes levando à marginalização de perspectivas e estilos de vida alternativos\footnote{Adaptado de Oxford Review\cite{oxfordreview}.}.

Ou seja, a classe dominante produz a cultura hegemônica (dominante) que promove sua própria concepção de mundo. Esta cultura abrange, expande e modifica conforme necessário um conjunto de signos, valores e práticas que justificam e sustentam a manutenção desta mesma classe no poder. Não foi ao acaso que Marx, na Ideologia Alemã declarou que “as ideias da classe dominante são as ideias dominantes”\cite{marx2007ideologia}. Esta dominação é construída através da dimensão cultural da sociedade. 

Assim sendo, as próprias relações de produção que ocupam uma posição especial na estrutura social não são exteriores a esta estrutura social e cultural que a justifica. Uma das grandes contribuições de Marx para o debate da economia política foi reconhecer o aspecto social do trabalho. O próprio “valor” conforme exposto por Marx, só existe como uma relação de troca, ou seja, uma relação social, e não como uma propriedade intrínseca do bem produzido.

Produzir uma “economia comunista” no sentido de elaborar apenas um plano técnico de produção levando em conta a tecnologia, matéria prima e ferramentas disponíveis, considerando que todo mundo já concorda com o comunismo, não é o maior desafio, é uma questão meramente técnica. Podemos lembrar das propostas apresentadas em “Towards a new Socialism”\cite{cockshott1993towards} ou “Ciber-comunismo: planificación económica, computadoras y democracia”\cite{cockshott2020cibercomunismo}, ambos do mesmo autor.

Não que seja simples, mas é uma questão técnica mais fácil de se trabalhar. A ciência e tecnologia necessárias já existem. O desafio que ainda resta é sobre como produzir uma cultura no qual as pessoas estejam comprometidas em construir e manter esse modo de produção comunista? Líderes comunistas históricos sentiram na prática a necessidade de abordar o aspecto cultural na superação do capitalismo, e eu gostaria de comentar rapidamente antes de concluir a sessão.

Podemos começar lembrando do fato de que a China realizou uma revolução cultural onde visava literalmente “revolucionar” a cultura, construindo e consolidando uma “nova cultura proletária” em detrimento de uma “antiga cultura burguesa”. Discute-se o sucesso desse esforço e os métodos que foram empregados, mas dificilmente comunistas discordam da necessidade de realizar uma revolução cultural. É possível que um dos problemas da URSS que levou ao fim do bloco soviético tenha sido exatamente a ausência de uma revolução cultural mais contundente. Ou seja, um dos possíveis motivos para a dissolução do bloco soviético pode consistir exatamente no aspecto cultural.

Não podemos ver a cultura como um mero reflexo, mas também é mais que um mero acessório do modo de produção, uma vez que aparentemente, sem uma cultura hegemônica adequada, um modo de produção pode falhar em se estabelecer a longo prazo. A experiência tem nos mostrado que não é possível esperar realizar primeiro uma revolução política e econômica, para apenas futuramente realizar uma revolução cultural de forma separada, todas estas dimensões (política, econômica e cultural) precisam ser modificadas de forma conjunta.

E isto não é, de fato, algo que passou despercebido pelo movimento comunista. Ho Chi Minh por exemplo, fala da necessidade dos membros do partido abandonarem o individualismo, que estes deveriam servir como modelos para os trabalhadores\cite{hochiminh1958moralidade}, isto é, na prática, a construção de uma nova cultura de valores coletivos em detrimento de uma cultura individualista. Em outro episódio, em um discurso intitulado “Construindo um Povo Socialista”, Ho Chi Minh explicou que um dos principais objetivos da revolução social vietnamita é construir um povo socialista com valores e costumes socialistas\cite{historical}.

Enver Hoxha vai falar da necessidade de reforçar a “consciência comunista” das massas trabalhadoras\cite{hoxha1974educacao}, o que é isto se não a construção de uma nova cultura? Che Guevara vai dizer que precisamos criar o Homem do século XXI\cite{guevara1965socialismo}. Kim Il Sungvai discutir sobre a necessidade de educar o povo e inspirar neste povo um amor ardente ao seu lugar de origem\cite{kim1955juche}, ideia que está nos pilares da Ideia Juche. Certamente poderíamos seguir citando outros revolucionários que caminham sob a mesma linha, mas para concluir, quero deixar apenas uma última citação:
\begin{quote}
Nos últimos vinte anos, para onde quer que essa nova força cultural tenha levado os seus ataques, verificou-se uma grande revolução, tanto no conteúdo ideológico como na forma (por exemplo na linguagem escrita). A sua influência é tão grande e o seu impacto tão poderoso que se afigura invencível onde quer que chega. A mobilização a que se procedeu ultrapassou a de qualquer outro período na China. Lu Sun foi o maior e mais corajoso porta-bandeira dessa nova força cultural. Comandante em chefe da revolução cultural chinesa, ele não foi apenas um grande homem de letras, foi também um grande pensador e um grande revolucionário. Lu Sun foi um homem de coluna vertebral tesa, sem sombra de servilismo nem obsequiosidade --- qualidade inestimável dos povos das colónias e semi-colónias. Na frente cultural, Lu Sun foi o representante da grande maioria da nação, o herói nacional mais correto, mais bravo, mais firme, mais fiel e mais ardente, um herói que abriu brecha e arrasou a cidadela do inimigo, um herói sem igual. A via de Lu Sun foi justamente a via da nova cultura nacional da China\cite{mao1940democracia}.
\end{quote}
Como exatamente isso deve ser feito é uma questão em aberto, porém não há discussão sobre que isso precisa ser feito. Trabalhando com um exemplo nacional, devemos lembrar do Augusto Boal e da sua posição de que: ``Pode ser que o teatro não seja revolucionário em si mesmo, mas estas formas teatrais são certamente um ensaio da revolução\cite{boal1991teatro}.'' Brecht também é conhecido por defender uma posição onde vê no teatro, e na arte no geral, uma ferramenta política a ser utilizada pelo movimento comunista. Não é apenas sobre construir a tecnologia necessária para que seja possível um novo modo de produção e distribuição de bens, mas também sobre a construção de novos valores e afetos, de uma nova concepção de mundo comunista, e consequentemente, de um novo modo de vida comunista.
\begin{quote}
Os reacionários -- especialmente os fascistas -- procuram evitar que essa situação seja devidamente compreendida, procuram escamotear a significação política global da arte e até procuram estetizar a política\ldots{} Essa estetização da política (com o embelezamento fascista da guerra) exige uma réplica: “A resposta do comunismo é a de politizar a arte\cite{benjamin1936obra}”\ldots{} o problema estético surge quando eles se perguntam como a arte deve se politizar\cite{konder2013marxistas}.
\end{quote}

\subsection{Economia Marxista 0}

A terceira componente é a economia marxista. É preciso investigar os aspectos econômicos da sociedade capitalista para saber como superá-lo. Podemos notar como na verdade o materialismo dialético, o socialismo científico e a economia marxista não podem ser entendidos como disciplinas isoladas. É preciso mobilizar todos os conhecimentos para uma efetiva transformação da realidade. Quanto a primeira fase do comunismo, Marx nos fornece alguns detalhes da organização socioeconômica que ele visava para esta etapa histórica:
\begin{quote}
O que o produtor deu à sociedade constitui sua cota individual de trabalho. Assim, por exemplo, a jornada social de trabalho compõe-se da soma das horas de trabalho Individual; o tempo Individual de trabalho de cada produtor em separado é a parte da jornada social de trabalho com que ele contribui, é sua participação nela. A sociedade entrega-lhe um bônus consignando que prestou tal ou qual quantidade de trabalho \ldots{} e com este bônus ele retira dos depósitos sociais de meios de consumo a parte equivalente à quantidade de trabalho que prestou. A mesma quantidade de trabalho que deu à sociedade sob uma forma, recebe-a desta sob uma outra forma diferente\ldots{} Variaram a forma e o conteúdo, porque sob as novas condições ninguém pode dar senão seu trabalho, e porque, de outra parte, agora nada pode passar a ser propriedade do indivíduo, fora dos meios individuais de consumo\cite{marx1875gotha}.
\end{quote}
Agora para entendermos melhor como se daria a construção desta sociedade, sobre como este bônus deve operar e como ele difere do dinheiro, acabamos tendo que nos dedicarmos ao estudo da economia marxista.
\begin{quote}
Quando Marx afirma que os certificados de trabalho não valem mais do que um ingresso de teatro, podemos inferir certas implicações:

(1) Os certificados não circulam; eles só podem ser trocados diretamente por bens de consumo.

(2) \ldots{} Somente a pessoa que realizou o trabalho poderia usá-los.

(3) Eles seriam cancelados após um único uso\ldots{} A loja, como uma organização comunitária, não precisa comprar mercadorias, apenas as recebe, portanto, seu único interesse nos comprovantes de trabalho é para fins de registro.

(4) Eles não serviriam como reserva de valor. Eles poderiam ter uma data de validade\cite{cockshott1993towards}\ldots{}
\end{quote}
Evidentemente, a forma com que esta ciência será abordada, irá depender de como entendemos o materialismo dialético, e isto impactará por sua vez no desenvolvimento do socialismo científico. Ou seja, apesar de tradicionalmente discutirmos que o marxismo tem três componentes, estes três componentes não podem ser entendidos de maneira isolada. Faço esta discussão com o propósito de contextualizar a defesa de que quando vamos estudar economia política, temos hoje ferramentas disponíveis mais interessantes e úteis do que as que Marx dispôs em sua época, de forma que é natural adotarmos novas abordagens. Em outras palavras, reafirmo a ideia expressa anteriormente de que a dialética hegeliana não é mais a principal ferramenta para realizar esta investigação.

Seguindo esta linha de raciocínio, só para citar alguns exemplos, Farjoun e Machover vão desenvolver uma formulação probabilística da economia política e da teoria do valor ao longo de suas principais obras\cite{farjoun2020laws,farjoun2022labor}. Paul Cockshott vai escrever uma contribuição para a continuação da tradição marxista que se inicia com o trabalho de Marx, mas que não abre mão do uso de ferramentas científicas modernas\cite{cockshott2012classical}.
\begin{quote}
Após uma análise das trocas básicas de mercadorias, passamos a examinar o dinheiro e o que Marx denominou ‘a forma do valor’. A obra de Marx permanece interessante, pois tentou aplicar uma análise formal à troca e ao dinheiro, utilizando as ferramentas então disponíveis: a lógica hegeliana. Tentamos reanalisar o aspecto formal do valor de troca utilizando uma tríade de conceitos extraídos de teorias modernas de sistemas formais:

(1) o conceito de espaço métrico,

(2) a noção de simetria e

(3) a ideia de assinatura\cite{cockshott2012classical}.
\end{quote}
Nesta mesma obra, Ian Wright vai também contribuir com abordagens probabilísticas e estatísticas que investigam temas como a arquitetura social do capitalismo\cite{Wright_2005,wright2008implicit} e a lei do valor\cite{wright2008emergence}. De certa forma, eu entendo por economia marxista simplesmente uma abordagem da economia que respeita os pressupostos do materialismo dialético. De modo geral podemos falar em:
\begin{itemize}
\item \textbf{Economia Marxista}: quando realizamos uma investigação econômica coerente com o materialismo dialético e nasce da crítica à economia política inglesa.
\item \textbf{Socialismo científico}: quando realizamos uma investigação e prática acerca da superação do capitalismo coerente com o materialismo dialético e nasce da crítica ao socialismo utópico francês, necessariamente coerente com a economia marxista.
\item \textbf{Materialismo Histórico}: quando realizamos uma investigação acerca da realidade histórica e social da humanidade em particular, e evidentemente, coerente com o materialismo dialético assim quanto a economia marxista.
\end{itemize}
O argumento de que o socialismo seria menos produtivo que o capitalismo é frequentemente mobilizado por críticos do socialismo. Embora eu não veja razões para que um modo de produção socialista seja intrinsecamente menos produtivo, afinal, ele pode se apropriar de todas as tecnologias existentes e utilizá-las de forma ainda mais racional, livre das restrições impostas pela busca do lucro, há um problema real que merece atenção: o ritmo de desenvolvimento das forças produtivas.

Diversos autores marxistas reconhecem que um dos diferenciais históricos do capitalismo em relação aos modos de produção anteriores foi o incentivo permanente à maximização da produtividade. Trata-se de um mérito histórico específico: nenhum modo de produção anterior estimulou de forma tão contínua o desenvolvimento das forças produtivas. Esse incentivo, no entanto, não é voluntário, mas coercitivo. O capitalista é compelido a priorizar o lucro sob pena de falência, e essa lógica o obriga a investir continuamente em produtividade de forma mais intensa que em qualquer outro modo de produção anterior da humanidade.

Nos modos de produção pré-capitalistas, a exploração operava predominantemente por vias extra econômicas. O aumento da exploração ocorria por meios como a ampliação de tributos ou a conquista de novas terras, isto é, a intensificação da coerção política direta. Mesmo quando havia pressão por maior produção, essa pressão tendia a se estabilizar após certo ponto. No capitalismo, ao contrário, a coerção econômica é permanente, forçando o reinvestimento contínuo.

Além disso, a ameaça constante do desemprego cria também para o trabalhador um incentivo coercitivo à elevação da produtividade. Diferentemente do servo, que não podia ser facilmente expulso da terra, o trabalhador assalariado pode ser substituído, o que intensifica a disciplina do trabalho. O problema que se coloca é que, no socialismo, essa coerção econômica capitalista é suprimida. Isso não implica menor produtividade em termos absolutos, mas pode implicar um ritmo menor de inovação tecnológica, caso não existam mecanismos substitutivos capazes de sustentar um desenvolvimento acelerado das forças produtivas. A China por exemplo, reintegrou com mecanismos de coerção econômica baseados no mercado, mantendo a lei do valor em funcionamento. Isso permite caracterizar a China menos como um modo de produção socialista e mais como uma forma específica de ditadura do proletariado operando sobre um modo de produção capitalista. 

Devemos entender que no vocabulário marxista, ditadura do proletarido não é sinônimo da ditadura militar da mesma forma que normalmente é utilizado pela mídia. No marxismo, devido ao caráter de classe do estado, se considera que enquanto existir classes, sempre existirá uma ditadura, pois a classe que controla o estado, exerce o monópolio da violência contra as outras classes, o que caracteriza a ``ditadura''. No capitalismo temos a ditadura da burguesia, na transição para o comunismo teríamos a ditadura do proletariado. Socialismo é outro termo que frequqnetemente aparece para descrever o período de transição para o comunismo, de fato, há um intenso debate sobre o que seria socialismo. Enquanto nunca existiu uma etapa chamada de socialismo nas obras de Marx, Lenin tem diferentes conceitos de socialismo ao longo da sua vida. Eu defendo a definição no qual um ``governo socialista'' não implica necessariamente em ter um ``modo de produção socialista''.
\begin{quote}
Creio que ninguém, ao analisar a questão do sistema econômico da Rússia, negou seu caráter de transição. Creio também que nenhum comunista negou que o termo “República Socialista Soviética” implica a determinação do poder soviético em realizar a transição para o socialismo, e não que o novo sistema econômico seja reconhecido como uma ordem socialista\cite{fleming1989lenin}.

Esta é uma posição semelhante a adotada pelo partido comunista do Vietnã.

A orientação socialista implica avançar em direção ao socialismo, ou, em outras palavras, estar em um período de transição para o socialismo. Quando, em 1921, surgiram dúvidas sobre o nome do país --- República Socialista Soviética ---, V. I. Lenin explicou que o termo significava que o governo soviético estava determinado a avançar para o socialismo, mas não significava, de forma alguma, que reconhecesse o novo regime econômico como socialismo. Da mesma forma, o termo “República Socialista do Vietnã” expressa apenas a determinação do Partido, do Estado e do povo do Vietnã de avançar para o socialismo, mas não significa que o Vietnã tenha concluído a construção do socialismo. A situação é semelhante na China\cite{tung2024socialist}. 
\end{quote}
Assim sendo, adoto uma definição intencionalmente ampla de socialismo como um governo que permita que os trabalhadores exerçam poder político e que esteja comprometido com a superação do modo de produção capitalista. Para fins de esclarecimento, podemos adotar a ditadura do proletariado como um período mais amplo conforme definido por Kim Il Sung:
\begin{quote}
Mas, aqui levanta-se uma outra questão: logo que se realize o comunismo num país ou numa região, quando o capitalismo ainda subsiste no mundo, que acontecerá à ditadura do proletariado? Enquanto a revolução mundial ainda não estiver acabada e o capitalismo e o imperialismo subsistirem, mesmo que o comunismo esteja realizado num país ou numa região, esta sociedade não poderá evitar a ameaça do imperialismo, nem a resistência dos inimigos do interior ligados aos inimigos do exterior. Em tais condições, o Estado não poderá extinguir-se e a ditadura do proletariado deverá permanecer sempre, mesmo na fase superior do comunismo. No caso da revolução estalar sucessivamente em todos os países do mundo e no caso do capitalismo se arruinar, triunfando a revolução socialista à escala mundial, o período de transição e a ditadura do proletariado coincidirão e, quando o período de transição chegar ao fim, a ditadura do proletariado já não será necessária e extinguir-se-ão as funções do Estado\cite{kim1967transicao}.
\end{quote}
Ou seja, a partir do momento em que os trabalhadores tomam o controle do Estado e passam a utilizá-lo para impor seus interesses históricos, tendo como objetivo a superação do modo de produção capitalista, inicia-se aquilo que designo como “socialismo”: um período de transição entre o capitalismo e o comunismo. Com a superação das relações capitalistas de produção, ingressa-se então no que Marx denomina de “primeira fase do comunismo”, caracterizada pela abolição do trabalho assalariado e pela distribuição da produção de acordo com a contribuição de cada trabalhador. Durante todo esse processo, enquanto persistirem classes sociais antagônicas, isto é, enquanto existir uma burguesia buscando retomar o poder, tanto interna quanto externamente, uma vez que o capitalismo constitui um sistema global, permanece também aquilo que se convencionou chamar de “ditadura do proletariado”: o controle do Estado pelos trabalhadores como instrumento de repressão da burguesia e de transformação das relações sociais. Nesse sentido, o definhamento do Estado, tal como proposto por Marx e Engels, pressupõe a abolição das classes sociais em escala internacional.

Além da alternativa chinesa, há outras alternativas, como a generalização de cooperativas, oficialmente defendida pela Coreia Popular\cite{kim1975construcao}. Porém esta alternativa também não elimina completamente a coerção econômica, já que a lei do valor persiste nas relações entre unidades produtivas. É preciso avalir se tais experiências demonstraram capacidade de sustentar um ritmo de desenvolvimento comparável ao das potêncis capitalistas, especialmente no plano militar.

Antes de concluir essa sessão, ainda que seja introdutóri, é necessário entendermos o que significa superar o capitalismo? De forma resumida, é superar a lei do valor. Dessa forma um governo socialista deve buscar construir as condições necessárias para que a lei do valor seja abolida. Por isso precisamos ter uma compreensão básica sobre o que é a lei do valor. Eu gostaria de discutir em que situações o valor corresponde ao preço.
\begin{quote}
Para que os preços a que as mercadorias são trocadas correspondam aproximadamente aos seus valores, nada mais é necessário do que 1) que a troca das várias mercadorias deixe de ser puramente acidental ou apenas ocasional; 2) no que diz respeito à troca direta de mercadorias, que essas mercadorias sejam produzidas em ambos os lados em quantidades aproximadamente suficientes para atender às necessidades mútuas, algo aprendido com a experiência mútua no comércio e, portanto, uma consequência natural do comércio contínuo; e 3) no que diz respeito à venda, que não haja nenhum monopólio natural ou artificial que permita a qualquer uma das partes contratantes vender mercadorias acima de seu valor ou obrigá-las a vender abaixo do valor. Por monopólio acidental entendemos um monopólio que um comprador ou vendedor adquire por meio de uma situação acidental de oferta ou demanda. A suposição de que as mercadorias das várias esferas de produção são vendidas pelo seu valor implica, naturalmente, que o seu valor é o centro de gravidade em torno do qual os seus preços flutuam, e que as suas subidas e descidas contínuas tendem a se equalizar\cite{marx2017capital3}.
\end{quote}
Esta forma do Marx abordar a questão, conduz diretamente a uma interpretação probabilística da lei do valor. A lei do valor nos diz que, em equilíbrio, o preço médio de uma mercadoria, a expressão monetária desta lei, tende a se correlacionar com seu valor, expresso quantitativamente através do número médio de horas necessárias para sua produção. Ou seja, se a sociedade produzir uma determinada mercadoria em quantidade suficiente para atender à demanda, o preço médio dessa mercadoria estará relacionado ao tempo médio necessário para sua produção. 

Comparando com outras mercadorias em equilíbrio, quanto maior o tempo de produção, maior o preço. Evidentemente, uma unidade individual de uma mercadoria pode ser vendida acima ou abaixo desse preço. E, como se trata de uma média, também é necessário que a mercadoria seja produzida em grande escala para que a lei do valor se manifeste dada a sua natureza estatística. Evidentemente, não quer dizer que o valor em si seja estatístico, mas sim que a maneira como o valor se manifesta nos preços só pode ser compreendida como uma tendência estatística: ele se torna visível nas médias quando observada toda a produção em grande escala ao longo do tempo. Este é um tema longe de ser consensual.

Portanto, a lei do valor não nos diz nada sobre o preço pelo qual a camisa autografada por Michael Jackson foi vendida em um leilão, pois este trata-se de um bem único, não de uma mercadoria produzida em grande escala. Assim, é importante notar que o preço serve como um sinal da relação entre demanda e oferta. Por exemplo, se uma mercadoria não está sendo produzida em quantidade suficiente para atender à demanda, isso causará uma distorção de preço, fazendo com que ela seja vendida muito acima de seu preço médio de equilíbrio. Por outro lado, a superprodução levará à venda abaixo da média.

É nítido então que o marxismo não ignora a distorção nos preços causado pela relação entre oferta e demanda de uma mercadoria, apenas a insere em um contexto mais amplo. É aqui também que surge as maiores críticas ao comunismo, pois defensores do capitalismo alegam que o livre-mercado é a melhor ferramenta que temos para processar as informações sobre que mercadorias a sociedade demanda e que recursos temos para produzir essas mercadorias, ou seja, para gerenciar a produção de bens.

É importante notar que a demanda não representa o interesse genuíno da sociedade por um determinado bem, apenas a demanda da fração da sociedade que tem condições de pagar por este bem. Além disso, os avanços tecnológicos que foram realizados pelo próprio capitalismo nos últimos 2 séculos tornam esta questão um tanto ultrapassada. Em um cenário de baixa tecnologia, o livre mercado é uma ferramenta rudimentar mas eficiente para estimar a demanda. Ele demanda a simples observação na variação do preços para estimar a relação de oferta e demanda no mercado, ou até mesmo ir ajustando empiricamente o preço das minhas mercadorias conforme a relação entre oferta e demanda que encontro. Porém hoje em dia há diversos recursos que podem ser utilizados para entender a demanda real por uma dada mercadoria na sociedade. Seria inocência acreditar que as maiores empresas capitalistas dependem apenas de observar a relação entre demanda e oferta para organizar sua produção e precificar seus produtos.

Outra maneira de entender a lei do valor é pensar nos produtos digitais. Para que o preço médio de uma mercadoria tenda para o seu valor, conforme descrito por Marx, é necessária a concorrência. Quando a concorrência é restrita, como no caso do monopólio, a expressão monetária da lei do valor fica distorcida. Um jogo digital, uma vez produzido, pode ser copiado infinitamente a custo virtualmente zero. Naturalmente, se a empresa A produz um jogo e o vende a qualquer preço inicial, a empresa B poder copiá-lo e vendê-lo mais barato, ou até o próprio usuáro poderia copiá-lo e distribuí-lo gratuitamente, então nesse cenário em que a lei do valor poderia operar sem barreiras, o preço tenderia a zero. 

No entanto, isso não acontece na realidade. Longe de negar a lei do valor, o que acontece é que as empresas de jogos impedem legalmente esse cenário por meio da propriedade intelectual. São criados então obstáculos que impedem outras empresas de vender o jogo, o que garante um monopólio e distorce a lei do valor.

Antes de concluir, é preciso relembrar que o desenvolvimento das forças produtivas leva a uma mudança nas relações de produção uma vez que elas se tornam uma restrição. Foi durante o capitalismo que desenvolvemos as forças produtivas a um nível em que podemos desenvolver softwares complexos, alcançando níveis sem precedentes de automação e integração de dados na história da humanidade. Estamos discutindo questões como a Indústria 4.0 e até mesmo a 5.0 em alguns países. Tamanho nível de desenvolvimento tecnologico tem o potencial de proporcionr uma vida mais confortável há toda humanidade, mas isso não interessa à classe que detém o poder político; ela usa sua posição privilegiada nas relações de produção para criar restrições burocráticas que impedem o desenvolvimento e o uso dessas ferramentas e, consequentemente, afetam negativamente a produção.

É importante ter aqui claramente a concepção de que ser comunista não é nem negar os avanços que o capitalismo trouxe, nem buscar retroceder na história. Pelo contrário, é acreditar que todo o progresso proporciondo pelo capitalismo pode ser amplificado com as relações sociais de produção adequadas. A discussão sobre a automação do processo de produção e sua relação com a teoria do valor não é novidade.
\begin{quote}
Uma terceira e última prova da veracidade da teoria do valor-trabalho é a prova por redução ao absurdo. Além disso, é a mais elegante e “moderna” das provas.

Imagine por um momento uma sociedade na qual o trabalho humano vivo desapareceu completamente, ou seja, uma sociedade na qual toda a produção foi 100\% automatizada... Mas imaginemos que esse desenvolvimento foi levado ao extremo e o trabalho humano foi completamente eliminado de todas as formas de produção e serviços. O valor pode continuar a existir nessas condições? Pode haver uma sociedade em que ninguém tenha renda, mas as mercadorias continuem a ter valor e a ser vendidas? Obviamente, tal situação seria absurda. Uma enorme quantidade de produtos seria produzida sem que essa produção gerasse qualquer renda, uma vez que nenhum ser humano estaria envolvido nessa produção. Mas alguém iria querer “vender” esses produtos para os quais não havia mais compradores!

É óbvio que a distribuição de produtos em tal sociedade não seria mais efetuada na forma de venda de mercadorias e, na verdade, a venda se tornaria ainda mais absurda devido à abundância produzida pela automação geral. Em outras palavras, uma sociedade na qual o trabalho humano fosse totalmente eliminado da produção, no sentido mais geral do termo, incluindo os serviços, seria uma sociedade na qual o valor de troca também teria sido eliminado. Isso prova a validade da teoria, pois, no momento em que o trabalho humano desaparece da produção, o valor também desaparece com ele\cite{mandel1967tratado}. 
\end{quote}
É também compartilhando a partir desta linha de raciocínio que Coggiola declara que a semente da superação do capitalismo está contida no próprio capitalismo. O aumento do capital constante (custos como instalações fixas, máquinas, terrenos, prédios, matérias primas, etc) em razão do capital variável (os custos com a mão de obra) devido aos constantes desenvolvimento tecnológicos empurraria a humanidade para uma situação propícia para a superação do capitalismo\cite{coggiola2024teoria}.

Erik Wright também faz uma discussão parecida apontando sobre a precarização do trabalho e do salário de grande parte dos trabalhadores e como isso poderia forçar a instituição de uma renda básica universal como forma de conter a crise\cite{wright2019anticapitalista}. Independente da estratégia adotada pela sociedade para lidar com as consequências do desenvolvimento da automação, é evidente que o desenvolvimento das forças produtivas no capitalismo caminham em uma direção de se chocar cada vez de forma mais violenta com as próprias relações de produção que definem o capitalismo, uma vez que o aumento de produtividade do capitalismo faz com que cada vez seja preciso menos trabalho para produzir as condições materiais necessárias para a manutenção da sociedade.

Em uma sociedade de produção abundante, o capitalismo não faz sentido. É por isso que o capitalismo precisa produzir escassez e novas necessidades. Neste contexto de abundância, o capitalismo perde o seu grande “mérito” de ser o mais eficiente motor que a humanidade já construiu para maximizar a produtividade, uma vez que não há mais esta necessidade. Para aumentar sua própria sobrevivência, o capitalismo ainda adota diferentes estratégias, como a expansão das suas fronteiras espaciais quando possível, e quando não mais possível, a destruição da base material através de guerras.

Até aqui, eu discuti apenas a manifestação monetária do valor dado pelo preço da mercadoria, ou de forma mais geral, a sua manifestação quantitativa dada pela quantidade de horas socialmente necessárias na produção da mercadoria. Mas é importante ter claramente que o valor não se reduz a isso.
\begin{quote}
O valor de uso e o valor não são, como sustenta Bohm-Bawerk, duas diferentes propriedades das coisas. O contraste entre ambos é provocado pelo contraste entre o método das ciências naturais, que trata a mercadoria como uma coisa, e o método sociológico, que trata de relações sociais “fundidas com coisas" . " O valor de uso expressa uma relação natural entre uma coisa e um homem, a existência de coisas para o homem. Mas o valor de troca representa a existência social das coisas"\cite{rubin1980teoria}...
\end{quote}
Algo que não havia ainda sido comentado é o caráter social do valor. Isto é, o valor de uma mercadoria não é uma propriedade intrínseca da mercadoria. Não é como o peso de um bem material, que é uma propriedade física desse material. Mas uma mercadoria só adquire valor dentro de um contexto social adequado. Até mesmo sua expressão quantitativa como a quantidade socialmente de horas necessárias, exige considerar a sociedade no qual a produção está inserida e conecta os produtores entre si, uma vez que o valor não é resultado da produção isolada deste produtor. Também conecta o produtor a possíveis consumidores, afinal a relação entre oferta e demanda não vai afetar apenas a expressão monetária do valor, mas a própria existência do mesmo.

Se você plantar uma laranjeira apenas para consumo pessoal, não há valor de troca nesta laranja. Se ninguém ver uso nesta laranja, mesmo que você tenha produzido para troca, não há valor de uso. Em ambos os casos, o valor cessa de existir. É preciso ter de form clara que o que determina o capitalismo são as relações sociais de produção, que são em si relações sociais entre indivíduos. O próprio valor de uma mercadoria é fruto de uma relação social.

Em uma economia socialista podemos utilizar como métrica para organizar a produção a quantidade média de horas necessárias para que uma dada mercadoria seja produzida\cite{cockshott1993towards}, mas isto não é mais o valor expresso na teoria do valor marxista, pois isto diz respeito especificamente à uma relação que existe no modo de produção capitalista.

Esta natureza social do capitalismo é constantemente ocultada pela cultura hegemônica capitalista. Esquecemos que a sociedade se organiza através da relação entre pessoas, e passamos a pensar que a sociedade se manifesta pela relação entre pessoas e coisas, ou muitas vezes, até mesmo apenas pela relação entre coisas. Aquilo que tem uma natureza social, passa a ser vista como se tivesse uma existência autônoma e é naturalizado como algo imutável e eterno, desde o valor de uma mercadoria, até a própria existência do capitalismo.

Esta é uma discussão sobre algo que constantemente ganha o nome de fetiche. Marx em seus escritos mais tardios muito pouco usou este termo. Quando pensamos no nosso exemplo da laranja, pudemos observar que uma simples produção visando a troca no mercado conecta todos produtores e consumidores entre si diretamente, tanto para a existência do valor quanto para a determinação do valor seja na sua expressão monetária em preço na sua sua expressão quantitativa em horas socialmente necessárias. Podemos ainda pensar que de forma direta ou indireta o produtor também se relaciona aos fornecedores de insumos necessários para a produção de laranja e de outras frutas, assim como a demanda por outras frutas no mercado. Podemos seguir este raciocínio de forma indefinida, não é difícil perceber que em última análise toda a sociedade está conectada.

Esta relação tende a ser ocultada no capitalismo e a relação entre pessoas nos parece uma relação entre coisas. Que uma troca entre um saco de 1kg laranja por uma moeda de R\$1, é uma relação entre as moedas e as laranjas apenas. Que 1kg de laranja vale 1 real pelas propriedades físicas da laranja e das moedas. Mas isso oculta toda natureza social do sistema de produção e as relações sociais envolvidas neste processo.

Mas isso não é um mero acaso que ocorre no capitalismo, isto decorre da forma com que o capitalismo organiza a sociedade. Pois ainda que seja determinado pelas relações sociais, a posição que as pessoas ocupam nesta arquitetura social do capitalismo é determinada pelas coisas que as pessoas possuem. Uma pessoa se torna capitalista por possuir capital. Agora, diferente de outros modos de produção que muitas vezes a posição que os indivíduos ocupavam na sociedade era definida predominantemente por relações sanguíneas, agora sua posição é definida pelos bens em sua posse. Ou seja, qualquer um, independente de sua origem, que passe a possuir capital, passa a ocupar a posição de capitalista. Isto é, se por um lado há uma objetificação das relações sociais, e por outro lado há uma humanização das coisas.

Evidentemente na prática, a posse das coisas e a mobilidade social é largamente herdada de nossos antepassados, uma vez que como discutido anteriormente, o capitalismo é, antes de tudo, uma forma de organizar as relações sociais de produção. E isso nos permite conectar toda esta discussão da dimensão econômica da sociedade com sua dimensão cultural. Rubin discutindo a forma social do capitalismo é obrigado a falar do estabelecimento de uma nova cultura, ainda que não intencionalmente.
\begin{quote}
Esta objetivação, ou " reificação"{} das relações de produção entre as pessoas sob a forma social de coisas, dá ao sistema econômico maior durabilidade, estabilidade e regularidade. O resultado é a " cristalização"{} das relações de produção entre as pessoas.Quando o tipo de produção dominante era a produção artesanal, na qual o objetivo era a " manutenção"{} do artesão, este ainda se considerava um " mestre artesão"{} e considerava seu rendimento como a fonte de sua " manutenção" , mesmo quando expandia sua empresa e já tinha se tornado, em essência, um capitalista que vivia do trabalho assalariado de seus operários. Ele ainda não considerava seu rendimento como o " lucro"{} do capital, nem seus meios de produção como " capital"\cite{rubin1980teoria}.
\end{quote}
Isto é, tanto da objetificação das relações sociais quanto da humanização das coisas não emergem de uma única ou poucas trocas isoladas. De fato, o mercado e a prática de trocas de mercadorias, inclusive com o uso de uma mercadoria comum ou dinheiro existiu desde muito antes do estabelecimento do modo de produção capitalista. Estas práticas precisam ocorrer com certa frequência e intensidade, para que só então ocorra de fato uma “reificação” ou “cristalização” das relações de produção. É apenas neste cenário que se estabelece então uma nova cultura hegemônica onde novos signos, valores e práticas passam a dominar a sociedade, consequentemente organizando também um novo modo de vida.

Concluindo então, entendo que a cristalização de novas relações sociais de produção implica simultaneamente a constituição de uma cultura hegemônica, na medida em que tais relações só conseguem existir enquanto modos socialmente reconhecidos, reiterados e vividos de agir, perceber e se orientar no mundo por pelo menos uma fração significativa da população que detém o monopólio da força. Se lido corretamente, o que Rubin descreve por um lado é uma descrição materialista da formação de uma nova cultura hegemônica determinada pela base material da sociedade, e por outro lado a descrição da importância da dimensão cultural no estabelecimento a longo prazo deste novo modo de produção.\bibliographystyle{plain}
\bibliography{tese,Artigo1/references,Artigo2/referencias}

\end{document}
